%% 3-9 放物線と y=ax^2（第12回編集委員会：組 V3-5）
%% 流れ：はてな（机から飛び出すボール）→ やってみよう（写真の位置を読む）→ ① 2乗に比例 → ② グラフ（放物線）
%%   → ③ a の符号と大きさ → ④ 変化の割合 → ⑤ 動かす y=a(x-p)^2+q → ⑥ 平方完成で頂点 → ⑦ 範囲と最大・最小 → ⑧ はてなの答え
%% 3-8（展開・平方完成）は別の組の章。使うときは「（3-8）」と1行で思い出させる。
\makeatletter
% 式を、まん中に1行で置く  \bkthrjline{…}
\newcommand\bkthrjline[1]{\par\addvspace{3pt}{\centering #1\par}\addvspace{3pt}}
% 説明の段の見出し（丸数字＋太字の1行）  \bkthrjstep{番号}{見出し}
\newcommand\bkthrjstep[2]{\par\addvspace{6pt}\Needspace{4\baselineskip}\noindent
  {\bkcircled[bkOrange]{#1}\enspace\sffamily\bfseries\color{bkOrange!75!black}#2}\par\nobreak\vspace{2pt}}
% 本文の中の小見出し（それだけで1行。下に3行以上入らないときは次のページへ）  \bkthrjhead{見出し}
\newcommand\bkthrjhead[1]{\par\addvspace{4pt}\Needspace{5\baselineskip}\noindent\textbf{#1}\par\nobreak}
% ボール  \bkthrjball{(x,y)}{半径}
\newcommand\bkthrjball[2]{\shade[ball color=bkOrange] #1 circle (#2);}
% 机（横から見た図）  \bkthrjdesk{(右上のかど)}{はば}{高さ}
\newcommand\bkthrjdesk[3]{\begin{scope}[shift={#1}]
  \fill[bkStem!55, rounded corners=0.4mm] (-#2,-2.2) rectangle (0,0);
  \fill[bkStem!70] (-#2+2,-2.2) rectangle (-#2+4,-#3);
  \fill[bkStem!70] (-4,-2.2) rectangle (-2,-#3);
\end{scope}}
\makeatother

\bkchapter{3}{9}{\R{放物線}{ほう|ぶつ|せん}と $y=ax^2$}{\bklevelrange{2}{3}}{\R{放物線}{ほう|ぶつ|せん}と $y=ax^2$}
\bkmichishirube{2-13,2-15,3-5,3-8}{3-10,3-16,5-1,5-2}

\begin{goal}
\begin{itemize}[label={\color{\bkTT}●}, leftmargin=1.6\zw, itemsep=0pt, topsep=1pt]
\item \bklevel{2}\enspace \R{表}{ひょう}から「$y$ は $x$ の2\R{乗}{じょう}に\R{比例}{ひ|れい}する」を\R{見}{み}つけ、$y=ax^2$ の\R{式}{しき}とグラフ（\R{放物線}{ほう|ぶつ|せん}）に\R{表}{あらわ}せる。
\item \bklevel{2}\enspace $a$ の\R{符号}{ふ|ごう}や\R{大}{おお}きさとグラフの\R{形}{かたち}の\R{関係}{かん|けい}、\R{変化}{へん|か}の\R{割合}{わり|あい}が\R{一定}{いっ|てい}でないことを\R{説明}{せつ|めい}できる。
\item \bklevel{3}\enspace \R{平方完成}{へい|ほう|かん|せい}で\R{放物線}{ほう|ぶつ|せん}の\R{頂点}{ちょう|てん}を\R{求}{もと}め、\R{範囲}{はん|い}があるときの\R{最大}{さい|だい}・\R{最小}{さい|しょう}を\R{求}{もと}められる。
\end{itemize}
\end{goal}

\begin{hatena}
\noindent\begin{minipage}[c]{0.55\linewidth}
\R{机}{つくえ}の\R{上}{うえ}を\R{転}{ころ}がってきたボールが、はしから\R{飛}{と}び\R{出}{だ}して、\R{床}{ゆか}に\R{落}{お}ちました。0.1\R{秒}{びょう}ごとにとった\R{写真}{しゃ|しん}を\R{重}{かさ}ねると、\R{右}{みぎ}の\R{図}{ず}のようになりました。

ボールは、まっすぐ\R{下}{した}には\R{落}{お}ちません。ななめの\R{直線}{ちょく|せん}にもなりません。きれいに\R{曲}{ま}がった\R{線}{せん}をえがいています。

この\R{曲線}{きょく|せん}は、どんな\R{形}{かたち}なのでしょう。\R{式}{しき}で\R{書}{か}けるのでしょうか。
\end{minipage}\hfill
\begin{minipage}[c]{0.42\linewidth}\centering
\begin{tikzpicture}[x=0.9mm, y=0.9mm]
  % ボールの中心の通り道：横 8k mm、たて -2k^2 mm（1 m を 40 mm でかいた。0.1秒ごとに k=0〜4）
  \bkthrjdesk{(0,-2)}{14}{34}
  \fill[bkGray!35] (-17,-36.2) rectangle (40,-34);
  \draw[bkGray!90!black, line width=0.8pt] (-17,-34) -- (40,-34);
  \draw[bkBlue!70, line width=0.8pt, dash pattern=on 2pt off 1.5pt, domain=0:32.6, samples=40] plot (\x,{-\x*\x/32});
  \foreach \k in {0,...,4} {\bkthrjball{(8*\k,{-2*\k*\k})}{2}}
  \bkthrjball{(-11,0)}{2}
  \draw[-{Stealth[length=3pt]}, bkInk, line width=0.7pt] (-7.5,0) -- (-3.5,0);
  \foreach \k in {1,...,4} {\draw[bkGray, line width=0.4pt, densely dotted] (8*\k,4) -- (8*\k,{-2*\k*\k+2.8});}
  \node[font=\sffamily\scriptsize, text=bkInk, anchor=south] at (16,4.5) {0.1\R{秒}{びょう}ごと};
\end{tikzpicture}
\end{minipage}
\end{hatena}

\clearpage
%% ---------------- やってみよう・思い出そう ----------------
\begin{yattemiyou}[\R{写真}{しゃ|しん}から、ボールの\R{位置}{い|ち}を\R{読}{よ}もう]
はてな？の\R{写真}{しゃ|しん}に、\qty{10}{cm}ごとの\R{目}{め}もりの\R{方眼}{ほう|がん}をかさねて、ボールのまん\R{中}{なか}の\R{位置}{い|ち}を\R{読}{よ}みとりました。「\R{横}{よこ}」は\R{机}{つくえ}のはしから\R{右}{みぎ}へ\R{進}{すす}んだ\R{長}{なが}さ、「\R{下}{さ}がった\R{長}{なが}さ」は\R{机}{つくえ}のはしの\R{高}{たか}さから\R{下}{さ}がった\R{長}{なが}さです。
\begin{bkcenter}
\begin{bktable}{colspec={Q[l,m,wd=38mm] *{5}{Q[c,m,wd=12mm]}}, rowsep=1.5pt}
\R{飛}{と}び\R{出}{だ}してからの\R{時間}{じ|かん} $t$（s） & 0 & 0.1 & 0.2 & 0.3 & 0.4 \\
\R{横}{よこ} $x$（m） & 0 & 0.2 & 0.4 & 0.6 & 0.8 \\
\R{下}{さ}がった\R{長}{なが}さ $y$（m） & 0 & 0.05 & 0.20 & 0.45 & 0.80 \\
\end{bktable}
\end{bkcenter}
\begin{enumerate}[label={\bkcircled[bkGreen]{\arabic*}}, leftmargin=2.4\zw, itemsep=1pt, topsep=3pt]
\item \R{横}{よこ} $x$ は、0.1\R{秒}{びょう}ごとに\R{何}{なん}mずつふえているかな？
\item \R{時間}{じ|かん} $t$ が0.1\R{秒}{びょう}から0.2\R{秒}{びょう}へ2\R{倍}{ばい}、0.3\R{秒}{びょう}へ3\R{倍}{ばい}、0.4\R{秒}{びょう}へ4\R{倍}{ばい}になると、\R{下}{さ}がった\R{長}{なが}さ $y$ は\R{何倍}{なん|ばい}になるかな？
\item 2-13 の「\R{落}{お}ちる\R{距離}{きょ|り}」の\R{式}{しき} $5t^2$ に、$t=0.1$、0.2、0.3、0.4 を\R{入}{い}れて\R{計算}{けい|さん}し、\R{表}{ひょう}の $y$ とくらべよう（\R{電卓}{でん|たく}を\R{使}{つか}ってよい）。
\end{enumerate}
\end{yattemiyou}

\noindent ①の\R{横}{よこ}は、0.1\R{秒}{びょう}ごとに \qty{0.2}{m} ずつ、\R{同}{おな}じだけふえています。1\R{秒}{びょう}あたりにすると $0.2\div0.1=2$ で、\qty{2}{m/s} の\R{一定}{いっ|てい}の\R{速}{はや}さです。\R{横}{よこ}の\R{動}{うご}きは、\R{時間}{じ|かん}に\R{比例}{ひ|れい}しています（2-15）。

②の\R{下}{さ}がった\R{長}{なが}さは、0.05 の4\R{倍}{ばい}が0.20、9\R{倍}{ばい}が0.45、16\R{倍}{ばい}が0.80 です。\R{時間}{じ|かん}が2\R{倍}{ばい}で4\R{倍}{ばい}、3\R{倍}{ばい}で9\R{倍}{ばい}、4\R{倍}{ばい}で16\R{倍}{ばい}。2\R{倍}{ばい}・3\R{倍}{ばい}・4\R{倍}{ばい}にはなっていません。

③では、たとえば $5\times0.3^2=5\times0.09=0.45$ で、\R{表}{ひょう}とぴったり\R{合}{あ}います。\R{下}{さ}がった\R{長}{なが}さは、2-13 で\R{作}{つく}った「\R{落}{お}ちる\R{距離}{きょ|り}」$5t^2$ そのものです。

\R{横}{よこ}は「\R{比例}{ひ|れい}」、たては「2\R{乗}{じょう}に\R{比例}{ひ|れい}」。この2つを\R{組}{く}み\R{合}{あ}わせると、ボールの\R{通}{とお}り\R{道}{みち}の\R{曲線}{きょく|せん}が\R{決}{き}まります。この\R{章}{しょう}で、その\R{曲線}{きょく|せん}の\R{正体}{しょう|たい}をつきとめましょう。

\begin{omoidasou}
\bkrecall{2-15}{$y=ax^2$ の\R{形}{かたち}を「$y$ は $x$ の2\R{乗}{じょう}に\R{比例}{ひ|れい}する」という。$x$ が2\R{倍}{ばい}になると、$y$ は $2\times2=4$ \R{倍}{ばい}。}
\bkrecall{2-13}{\R{落}{お}ちる\R{距離}{きょ|り}は $5t^2$（m）。$t$ に $1+h$ を\R{入}{い}れるときは、かっこごと\R{入}{い}れて $5(1+h)^2$。}
\bkrecall{3-2}{\R{点}{てん} $(x,\ y)$ は、\R{原点}{げん|てん}から\R{右}{みぎ}へ $x$、\R{上}{うえ}へ $y$ の\R{所}{ところ}。マイナスなら\R{左}{ひだり}へ・\R{下}{した}へ。}
\bkrecall{3-5}{\R{変化}{へん|か}の\R{割合}{わり|あい}＝$y$ の\R{変化量}{へん|か|りょう}÷$x$ の\R{変化量}{へん|か|りょう}。1\R{次関数}{じ|かん|すう} $y=ax+b$ では、いつも $a$。}
\bkrecall{3-8}{$(a+b)^2=a^2+2ab+b^2$。\R{平方完成}{へい|ほう|かん|せい}：$x^2+6x+10=(x+3)^2+1$。}
\end{omoidasou}

\bkfukidashi{sora}{\R{横}{よこ}はまっすぐ、たては2\R{乗}{じょう}。2つを\R{合}{あ}わせると、どんな\R{線}{せん}になるんだろう？}

%% ---------------- 数学の道具 ----------------
\bksection{doguu}{\R{数学}{すう|がく}の\R{道具}{どう|ぐ}}

\bksubsection{1}{2\R{乗}{じょう}に\R{比例}{ひ|れい}する：$y=ax^2$}

\leavevmode\label{bk:3-9:nijo}やってみようの\R{下}{さ}がった\R{長}{なが}さ $y$ が、\R{時間}{じ|かん} $t$ とどんな\R{関係}{かん|けい}にあるか、\R{表}{ひょう}で\R{調}{しら}べます。\R{比例}{ひ|れい}なら、$y\div t$ が\R{一定}{いっ|てい}になるはずでした（2-15）。\R{計算}{けい|さん}してみると、\R{次}{つぎ}の\R{表}{ひょう}の4\R{行}{ぎょう}めのように、0.5、1、1.5、2 とふえていきます。\R{一定}{いっ|てい}ではないので、\R{比例}{ひ|れい}ではありません。

そこで、$t$ のかわりに $t^2$（$t\times t$）でわってみます（5\R{行}{ぎょう}め）。
\begin{bkcenter}
\begin{bktable}{colspec={Q[l,m,wd=40mm] *{4}{Q[c,m,wd=15mm]}}, rowsep=1.5pt}
\R{時間}{じ|かん} $t$（s） & 0.1 & 0.2 & 0.3 & 0.4 \\
$t^2$（$t\times t$） & 0.01 & 0.04 & 0.09 & 0.16 \\
\R{下}{さ}がった\R{長}{なが}さ $y$（m） & 0.05 & 0.20 & 0.45 & 0.80 \\
$y\div t$ & 0.5 & 1 & 1.5 & 2 \\
\SetCell{bg=bkYellow!30} $y\div t^2$ & 5 & 5 & 5 & 5 \\
\end{bktable}
\bktabcaption{$y\div t$ はふえていくが、$y\div t^2$ はいつも5。}
\end{bkcenter}
たとえば $t=0.3$ では、$0.45\div0.09=5$ です。$y\div t^2$ は、どの\R{時間}{じ|かん}でも5になりました。$y$ は $t^2$ に\R{比例}{ひ|れい}していて、$y=5t^2$ と\R{書}{か}けます。

このように、$y$ が $x^2$ に\R{比例}{ひ|れい}するとき、「$y$ は $x$ の\textbf{2\R{乗}{じょう}に\R{比例}{ひ|れい}する}」といい、\R{式}{しき}は
\bkthrjline{\fbox{$y=ax^2$}}
と\R{書}{か}けます。\bkidx{にじょうにひれい}{2\R{乗}{じょう}に\R{比例}{ひ|れい}}$a$ を\textbf{\R{比例定数}{ひ|れい|てい|すう}}といいます。\R{比例定数}{ひ|れい|てい|すう}は、$y=ax^2$ の\R{両辺}{りょう|へん}を $x^2$ でわって（2-14）、$a=\dfrac{y}{x^2}$ で\R{求}{もと}められます（$x$ が0のときはわれないので、0でない $x$ を\R{使}{つか}います）。\R{落}{お}ちる\R{距離}{きょ|り}の5は、\R{地球}{ち|きゅう}の\R{重力}{じゅう|りょく}で、1\R{秒}{びょう}ごとに\R{約}{やく} \qty{10}{m/s} ずつ\R{速}{はや}くなることから\R{出}{で}てきた\R{数}{かず}でした（2-13。くわしくは \qty{9.8}{m/s^2}、5-5）。

\bkthrjhead{なぜ、2\R{倍}{ばい}で4\R{倍}{ばい}になるの？　その1：\R{式}{しき}で}
$y=ax^2$ の $x$ を2\R{倍}{ばい}の $2x$ にして、かっこごと\R{入}{い}れます（2-13）。
\[ a(2x)^2=a\times2x\times2x=4\times ax^2 \]
もとの $ax^2$ の4\R{倍}{ばい}です。2が2\R{回}{かい}かけられるので、$2\times2=4$ \R{倍}{ばい}になるのです。3\R{倍}{ばい}なら $3\times3=9$ \R{倍}{ばい}、4\R{倍}{ばい}なら $4\times4=16$ \R{倍}{ばい}です。$a$ がどんな\R{数}{かず}でも\R{同}{おな}じです。

\bkthrjhead{なぜ、2\R{倍}{ばい}で4\R{倍}{ばい}になるの？　その2：\R{正方形}{せい|ほう|けい}の\R{面積}{めん|せき}で}
\noindent\begin{minipage}[c]{0.5\linewidth}
1\R{辺}{ぺん}が $x$ の\R{正方形}{せい|ほう|けい}の\R{面積}{めん|せき}は、$x\times x=x^2$ です。\R{面積}{めん|せき}は、1\R{辺}{ぺん}の2\R{乗}{じょう}に\R{比例}{ひ|れい}します（$y=x^2$、\R{比例定数}{ひ|れい|てい|すう}は1）。1\R{辺}{ぺん}を2\R{倍}{ばい}にすると、\R{右}{みぎ}の\R{図}{ず}のように、もとの\R{正方形}{せい|ほう|けい}がたてに2つ、\R{横}{よこ}に2つで、$2\times2=4$ こ\R{入}{はい}ります。1\R{辺}{ぺん}を3\R{倍}{ばい}にすると、$3\times3=9$ こです。
\end{minipage}\hfill
\begin{minipage}[c]{0.46\linewidth}\centering
\begin{tikzpicture}[x=1mm, y=1mm]
  \foreach \n/\xo in {1/0, 2/12, 3/32} {
    \foreach \i in {1,...,\n} {\foreach \j in {1,...,\n} {
      \filldraw[fill=bkBlue!18, draw=bkBlue!70!black, line width=0.5pt] ({\xo+(\i-1)*6},{(\j-1)*6}) rectangle ++(6,6);}}
    \pgfmathtruncatemacro\nn{\n*\n}
    \node[font=\sffamily\footnotesize, anchor=north] at ({\xo+\n*3},-1.8) {1\R{辺}{ぺん}$\times$\n};
    \node[font=\sffamily\footnotesize, anchor=north] at ({\xo+\n*3},-6.3) {\nn こ};}
\end{tikzpicture}
\end{minipage}

\bkthrjhead{\R{比例}{ひ|れい}とくらべる}
\R{同}{おな}じ $x$ の\R{値}{あたい}で、\R{比例}{ひ|れい} $y=2x$ と、2\R{乗}{じょう}に\R{比例}{ひ|れい} $y=2x^2$ をくらべてみましょう。
\begin{bkcenter}
\begin{bktable}{colspec={Q[l,m,wd=32mm] *{4}{Q[c,m,wd=14mm]}}, rowsep=1.5pt}
$x$ & 1 & 2 & 3 & 4 \\
$y=2x$（\R{比例}{ひ|れい}） & 2 & 4 & 6 & 8 \\
$y=2x^2$（2\R{乗}{じょう}に\R{比例}{ひ|れい}） & 2 & 8 & 18 & 32 \\
\end{bktable}
\bktabcaption{$x$ が1から4へ4\R{倍}{ばい}になると、$2x$ は4\R{倍}{ばい}、$2x^2$ は16\R{倍}{ばい}。}
\end{bkcenter}
$x=1$ では\R{同}{おな}じ2ですが、$x$ が\R{大}{おお}きくなるほど、2\R{乗}{じょう}に\R{比例}{ひ|れい}するほうがずっと\R{速}{はや}くふえます。$x=4$ では、$2\times4^2=2\times16=32$ です。

\bksubsection{2}{グラフは「\R{放物線}{ほう|ぶつ|せん}」}

\leavevmode\label{bk:3-9:graph}いちばんかんたんな $y=x^2$（\R{比例定数}{ひ|れい|てい|すう}1）のグラフをかきます。$x$ には、マイナスの\R{数}{かず}も\R{入}{い}れます。マイナスの\R{数}{かず}の2\R{乗}{じょう}は、$(-3)^2=(-3)\times(-3)=9$ のようにプラスです（2-2：マイナス×マイナスはプラス）。
\begin{bkcenter}
\begin{bktable}{colspec={Q[l,m,wd=14mm] *{7}{Q[c,m,wd=11mm]}}, rowsep=1.5pt}
$x$ & $-3$ & $-2$ & $-1$ & 0 & 1 & 2 & 3 \\
$y=x^2$ & 9 & 4 & 1 & 0 & 1 & 4 & 9 \\
\end{bktable}
\bktabcaption{$y=x^2$ の\R{表}{ひょう}。0をはさんで、\R{同}{おな}じ\R{数}{かず}が\R{左右}{さ|ゆう}にならぶ。}
\end{bkcenter}

\noindent\begin{minipage}[c]{0.47\linewidth}\centering
\begin{tikzpicture}
\begin{axis}[bkplot, width=56mm, height=58mm, xmin=-3.6, xmax=3.8, ymin=-0.8, ymax=10.2, xtick={-3,-2,-1,1,2,3}, ytick={2,4,6,8},
  xlabel={$x$}, ylabel={$y$}]
  \addplot[only marks, mark=*, mark size=1.8pt, bkInk] coordinates {(-3,9) (-2,4) (-1,1) (0,0) (1,1) (2,4) (3,9)};
  \addplot[only marks, mark=*, mark size=1.4pt, bkGreen!70!black] coordinates {(-2.5,6.25) (-1.5,2.25) (-0.5,0.25) (0.5,0.25) (1.5,2.25) (2.5,6.25)};
\end{axis}
\end{tikzpicture}
\bkfigcaption{\R{黒}{くろ}は\R{表}{ひょう}の\R{点}{てん}。\R{緑}{みどり}は $x=\pm0.5$、$\pm1.5$、$\pm2.5$ の\R{点}{てん}。}
\end{minipage}\hfill
\begin{minipage}[c]{0.47\linewidth}\centering
\begin{tikzpicture}
\begin{axis}[bkplot, width=56mm, height=58mm, xmin=-3.6, xmax=3.8, ymin=-0.8, ymax=10.2, xtick={-3,-2,-1,1,2,3}, ytick={2,4,6,8},
  xlabel={$x$}, ylabel={$y$}]
  \addplot[bkcurve, domain=-3.15:3.15] {x^2};
  \addplot[only marks, mark=*, mark size=1.5pt, bkInk] coordinates {(-3,9) (-2,4) (-1,1) (0,0) (1,1) (2,4) (3,9)};
  \node[font=\sffamily\footnotesize, text=bkBlue!80!black, anchor=west] at (axis cs:2.3,2.5) {$y=x^2$};
\end{axis}
\end{tikzpicture}
\bkfigcaption{\R{点}{てん}を、なめらかな\R{曲線}{きょく|せん}でつなぐ。}
\end{minipage}

\medskip
\R{表}{ひょう}の\R{点}{てん}（\R{黒}{くろ}）をかいたら、$x=0.5$、$1.5$ などのあいだの\R{点}{てん}（\R{緑}{みどり}）も\R{計算}{けい|さん}してかいてみます。たとえば $1.5^2=2.25$ です。\R{点}{てん}をふやすほど、\R{点}{てん}はなめらかな\R{曲線}{きょく|せん}の\R{上}{うえ}にならびます。

\bkthrjhead{なぜ、おれ\R{線}{せん}でつながないの？}
\noindent\begin{minipage}[c]{0.6\linewidth}
\R{点}{てん} $(1,\ 1)$ と $(2,\ 4)$ を\R{直線}{ちょく|せん}でつなぐと、まん\R{中}{なか}の $x=1.5$ では、\R{高}{たか}さが1と4のまん\R{中}{なか}の $(1+4)\div2=2.5$ になります。でも、\R{本当}{ほん|とう}の $y$ は $1.5^2=2.25$ で、\R{直線}{ちょく|せん}より\R{少}{すこ}し\R{下}{した}です。$y=x^2$ のグラフは、\R{点}{てん}と\R{点}{てん}のあいだで、\R{直線}{ちょく|せん}より\R{下}{した}をとおって\R{曲}{ま}がっています。だから、おれ\R{線}{せん}ではなく、なめらかな\R{曲線}{きょく|せん}でつなぎます。
\end{minipage}\hfill
\begin{minipage}[c]{0.36\linewidth}\centering
\begin{tikzpicture}
\begin{axis}[bkplot, width=42mm, height=44mm, xmin=0.6, xmax=2.3, ymin=0.4, ymax=4.6, xtick={1,1.5,2}, ytick={1,4},
  yticklabel style={/pgf/number format/fixed}, axis lines=left, xlabel={$x$}, ylabel={$y$}]
  \addplot[bkcurve, domain=0.8:2.15] {x^2};
  \addplot[bkGray, line width=0.9pt, dash pattern=on 2pt off 1.5pt] coordinates {(1,1) (2,4)};
  \addplot[only marks, mark=*, mark size=1.5pt, bkInk] coordinates {(1,1) (2,4)};
  \addplot[only marks, mark=*, mark size=1.5pt, bkGreen!70!black] coordinates {(1.5,2.25)};
  \addplot[only marks, mark=o, mark size=1.6pt, bkGray, line width=0.7pt] coordinates {(1.5,2.5)};
  \node[font=\sffamily\scriptsize, text=bkGray, anchor=south east, inner sep=1pt] at (axis cs:1.44,2.6) {2.5};
  \node[font=\sffamily\scriptsize, text=bkGreen!50!black, anchor=north west, inner sep=1pt] at (axis cs:1.56,2.1) {2.25};
\end{axis}
\end{tikzpicture}
\end{minipage}

\medskip
このような\R{形}{かたち}の\R{曲線}{きょく|せん}を、\textbf{\R{放物線}{ほう|ぶつ|せん}}といいます。\bkidx{ほうぶつせん}{\R{放物線}{ほう|ぶつ|せん}}「\R{物}{もの}を\R{放}{はな}ったときにえがく\R{線}{せん}」という\R{意味}{い|み}の\R{名前}{な|まえ}です。はてな？のボールの\R{通}{とお}り\R{道}{みち}も、じつは\R{放物線}{ほう|ぶつ|せん}です（\pageref{bk:3-9:kotae}ページで\R{確}{たし}かめます）。\R{英語}{えい|ご}では parabola（パラボラ）といいます。

\bkthrjhead{\R{放物線}{ほう|ぶつ|せん}の\R{形}{かたち}の\R{特}{とく}ちょう}
\noindent\begin{minipage}[c]{0.56\linewidth}
\begin{enumerate}[label={\bkcircled[bkOrange]{\arabic*}}, leftmargin=2.4\zw, itemsep=2pt, topsep=0pt]
\item \textbf{\R{左右対称}{さ|ゆう|たい|しょう}}\par\noindent $x=3$ と $x=-3$ は、どちらも $y=9$ です。2\R{乗}{じょう}すると、\R{符号}{ふ|ごう}のちがいが\R{消}{け}えるからです。だから、グラフを $y$ \R{軸}{じく}で\R{折}{お}ると、\R{左右}{さ|ゆう}がぴったり\R{重}{かさ}なります。このように、\R{折}{お}って\R{重}{かさ}なる\R{形}{かたち}を\textbf{\R{線対称}{せん|たい|しょう}}な\R{形}{かたち}といいます。
\item \textbf{\R{軸}{じく}}\par\noindent \R{折}{お}り\R{目}{め}になる\R{直線}{ちょく|せん}を、\R{放物線}{ほう|ぶつ|せん}の\textbf{\R{軸}{じく}}といいます。$y=x^2$ の\R{軸}{じく}は $y$ \R{軸}{じく}（$x=0$ の\R{直線}{ちょく|せん}）です。\bkidx{ほうぶつせんのじく}{\R{放物線}{ほう|ぶつ|せん}の\R{軸}{じく}}
\item \textbf{\R{頂点}{ちょう|てん}}\par\noindent \R{放物線}{ほう|ぶつ|せん}と\R{軸}{じく}の\R{交}{まじ}わる\R{点}{てん}を、\R{放物線}{ほう|ぶつ|せん}の\textbf{\R{頂点}{ちょう|てん}}といいます。$y=x^2$ の\R{頂点}{ちょう|てん}は\R{原点}{げん|てん} $(0,\ 0)$ で、グラフのいちばん\R{低}{ひく}い\R{点}{てん}です。\bkidx{ちょうてん}{\R{頂点}{ちょう|てん}}
\end{enumerate}
\end{minipage}\hfill
\begin{minipage}[c]{0.4\linewidth}\centering
\begin{tikzpicture}
\begin{axis}[bkplot, width=50mm, height=50mm, xmin=-3.3, xmax=3.6, ymin=-1.2, ymax=10, xtick={-3,3}, ytick=\empty,
  xlabel={$x$}, ylabel={$y$}]
  \addplot[bkcurve, domain=-3.1:3.1] {x^2};
  \addplot[bkRed, line width=1pt, dash pattern=on 3pt off 2pt] coordinates {(0,-1.1) (0,9.6)};
  \addplot[only marks, mark=*, mark size=2pt, bkRed] coordinates {(0,0)};
  \draw[{Stealth[length=3pt]}-{Stealth[length=3pt]}, bkGreen!70!black, line width=0.7pt] (axis cs:-2.8,7.84) -- (axis cs:2.8,7.84);
  \node[font=\sffamily\scriptsize, text=bkRed!80!black, anchor=north west] at (axis cs:0.2,-0.2) {\R{頂点}{ちょう|てん}};
  \node[font=\sffamily\scriptsize, text=bkRed!80!black, anchor=west] at (axis cs:0.25,5) {\R{軸}{じく}};
  \node[font=\sffamily\scriptsize, text=bkGreen!50!black, anchor=south] at (axis cs:-1.4,8) {\R{折}{お}ると\R{重}{かさ}なる};
\end{axis}
\end{tikzpicture}
\end{minipage}

\medskip
どんな\R{数}{かず}も、2\R{乗}{じょう}すると0か、0より\R{大}{おお}きくなります（2-8）。だから、$y=x^2$ のグラフは、\R{原点}{げん|てん}のほかは、ぜんぶ $x$ \R{軸}{じく}より\R{上}{うえ}にあります。

\begin{taisetsu}
\R{放物線}{ほう|ぶつ|せん}をかく\R{手順}{て|じゅん}\par
①\enspace \R{頂点}{ちょう|てん}と、その\R{左右}{さ|ゆう}の\R{点}{てん}を3つずつくらい\R{計算}{けい|さん}して、\R{表}{ひょう}にする。\par
②\enspace \R{点}{てん}を\R{方眼}{ほう|がん}にかく。\R{左右対称}{さ|ゆう|たい|しょう}になっているかを\R{確}{たし}かめる。\par
③\enspace \R{点}{てん}を、なめらかな\R{曲線}{きょく|せん}でつなぐ。\R{頂点}{ちょう|てん}のまわりは、とがらせずに\R{丸}{まる}くかく。\par
④\enspace \R{線}{せん}の\R{先}{さき}は、\R{点}{てん}の\R{外}{そと}まで\R{少}{すこ}しのばしてよい（\R{式}{しき}がそのまま\R{続}{つづ}くとき）。
\end{taisetsu}

\bksubsection{3}{$a$ の\R{符号}{ふ|ごう}と\R{大}{おお}きさで、\R{形}{かたち}が\R{変}{か}わる}

\leavevmode\label{bk:3-9:a}\R{比例定数}{ひ|れい|てい|すう} $a$ を\R{変}{か}えると、\R{放物線}{ほう|ぶつ|せん}の\R{形}{かたち}はどう\R{変}{か}わるでしょう。$a=2$、1、$\dfrac{1}{2}$ の3つで\R{表}{ひょう}を\R{作}{つく}ります。
\begin{bkcenter}
\begin{bktable}{colspec={Q[l,m,wd=24mm] *{5}{Q[c,m,wd=12mm]}}, rowsep=1.5pt}
$x$ & $-2$ & $-1$ & 0 & 1 & 2 \\
$y=2x^2$ & 8 & 2 & 0 & 2 & 8 \\
$y=x^2$ & 4 & 1 & 0 & 1 & 4 \\
$y=\dfrac{1}{2}x^2$ & 2 & 0.5 & 0 & 0.5 & 2 \\
\end{bktable}
\end{bkcenter}

\noindent\begin{minipage}[c]{0.47\linewidth}\centering
\begin{tikzpicture}
\begin{axis}[bkplot, width=56mm, height=56mm, xmin=-2.7, xmax=3, ymin=-0.6, ymax=8.8, xtick={-2,-1,1,2}, ytick={2,4,6,8},
  xlabel={$x$}, ylabel={$y$}]
  \addplot[bkRed, line width=1.6pt, domain=-2.05:2.05] {2*x^2};
  \addplot[bkcurve, domain=-2.6:2.6] {x^2};
  \addplot[bkGreen!70!black, line width=1.6pt, domain=-2.6:2.6] {0.5*x^2};
  \node[font=\sffamily\scriptsize, text=bkRed!80!black, anchor=south west] at (axis cs:2.1,8.5) {$y=2x^2$};
  \node[font=\sffamily\scriptsize, text=bkBlue!80!black, anchor=west] at (axis cs:2.62,6.6) {$y=x^2$};
  \node[font=\sffamily\scriptsize, text=bkGreen!50!black, anchor=west] at (axis cs:2.62,3.0) {$y=\frac{1}{2}x^2$};
\end{axis}
\end{tikzpicture}
\bkfigcaption{$a$ が\R{大}{おお}きいほど、\R{細}{ほそ}い\R{形}{かたち}になる。}
\end{minipage}\hfill
\begin{minipage}[c]{0.47\linewidth}
どれも、\R{原点}{げん|てん}が\R{頂点}{ちょう|てん}で、$y$ \R{軸}{じく}が\R{軸}{じく}です。$x=2$ のときを\R{見}{み}ると、$y$ は $2\times4=8$、$1\times4=4$、$\dfrac{1}{2}\times4=2$ です。\R{同}{おな}じ $x$ でくらべると、$y$ は $a$ に\R{比例}{ひ|れい}しています。$a$ が\R{大}{おお}きいほど、\R{同}{おな}じ\R{横}{よこ}の\R{位置}{い|ち}で\R{高}{たか}くなるので、\R{放物線}{ほう|ぶつ|せん}はたてに\R{引}{ひ}きのばされて、\R{細}{ほそ}い\R{形}{かたち}になります。$a$ が\R{小}{ちい}さいほど、\R{平}{ひら}たく\R{開}{ひら}いた\R{形}{かたち}になります。
\end{minipage}

\bkthrjhead{$a$ がマイナスのとき}
$y=-x^2$ の\R{表}{ひょう}を\R{作}{つく}ると、$y=x^2$ の\R{数}{かず}に、すべてマイナスがつきます。$x=2$ なら $-(2^2)=-4$ です（2\R{乗}{じょう}を\R{先}{さき}に\R{計算}{けい|さん}してから、マイナスをつけます。$(-2)^2$ とはちがいます。2-12）。
\begin{bkcenter}
\begin{bktable}{colspec={Q[l,m,wd=24mm] *{5}{Q[c,m,wd=12mm]}}, rowsep=1.5pt}
$x$ & $-2$ & $-1$ & 0 & 1 & 2 \\
$y=x^2$ & 4 & 1 & 0 & 1 & 4 \\
$y=-x^2$ & $-4$ & $-1$ & 0 & $-1$ & $-4$ \\
$y=-\dfrac{1}{2}x^2$ & $-2$ & $-0.5$ & 0 & $-0.5$ & $-2$ \\
\end{bktable}
\end{bkcenter}

\noindent\begin{minipage}[c]{0.47\linewidth}
どの\R{点}{てん}も、$y=x^2$ の\R{点}{てん}を $x$ \R{軸}{じく}で\R{折}{お}り\R{返}{かえ}した\R{所}{ところ}にきます。だから、$y=-x^2$ のグラフは、$y=x^2$ を\R{上下}{じょう|げ}にひっくり\R{返}{かえ}した\R{形}{かたち}になります。\R{頂点}{ちょう|てん}の\R{原点}{げん|てん}が、こんどはいちばん\R{高}{たか}い\R{点}{てん}です。

$a$ がプラスのときは、\R{谷}{たに}のような\textbf{\R{上}{うえ}に\R{開}{ひら}いた\R{形}{かたち}}、マイナスのときは、\R{山}{やま}のような\textbf{\R{下}{した}に\R{開}{ひら}いた\R{形}{かたち}}です。
\end{minipage}\hfill
\begin{minipage}[c]{0.47\linewidth}\centering
\begin{tikzpicture}
\begin{axis}[bkplot, width=56mm, height=56mm, xmin=-2.7, xmax=4.4, ymin=-4.8, ymax=4.8, xtick={-2,2}, ytick={-4,-2,2,4},
  xlabel={$x$}, ylabel={$y$}]
  \addplot[bkcurve, domain=-2.15:2.15] {x^2};
  \addplot[bkRed, line width=1.6pt, domain=-2.15:2.15] {-x^2};
  \addplot[bkGreen!70!black, line width=1.6pt, domain=-2.6:2.6] {-0.5*x^2};
  \node[font=\sffamily\scriptsize, text=bkBlue!80!black, anchor=west] at (axis cs:2.2,4.4) {$y=x^2$};
  \node[font=\sffamily\scriptsize, text=bkRed!80!black, anchor=west] at (axis cs:2.2,-4.4) {$y=-x^2$};
  \node[font=\sffamily\scriptsize, text=bkGreen!50!black, anchor=west] at (axis cs:2.65,-3.0) {$y=-\frac{1}{2}x^2$};
\end{axis}
\end{tikzpicture}
\bkfigcaption{$a$ がマイナスなら、\R{下}{した}に\R{開}{ひら}いた\R{形}{かたち}。}
\end{minipage}

\begin{taisetsu}
$y=ax^2$ のグラフは、\R{原点}{げん|てん}を\R{頂点}{ちょう|てん}とし、$y$ \R{軸}{じく}を\R{軸}{じく}とする\R{放物線}{ほう|ぶつ|せん}。\par\vspace{2pt}
$a>0$ なら\R{上}{うえ}に\R{開}{ひら}いた\R{形}{かたち}（\R{頂点}{ちょう|てん}がいちばん\R{低}{ひく}い）。$a<0$ なら\R{下}{した}に\R{開}{ひら}いた\R{形}{かたち}（\R{頂点}{ちょう|てん}がいちばん\R{高}{たか}い）。\par\vspace{2pt}
$a$ の\R{大}{おお}きさ（\R{符号}{ふ|ごう}をのぞいた\R{大}{おお}きさ）が\R{大}{おお}きいほど、\R{細}{ほそ}い\R{形}{かたち}になる。
\end{taisetsu}

\bksubsection{4}{\R{変化}{へん|か}の\R{割合}{わり|あい}は、\R{一定}{いっ|てい}ではない}

\leavevmode\label{bk:3-9:wariai}1\R{次関数}{じ|かん|すう} $y=ax+b$ の\R{変化}{へん|か}の\R{割合}{わり|あい}は、どこでくらべても $a$ で\R{一定}{いっ|てい}でした（3-5）。だからグラフは\R{直線}{ちょく|せん}です。$y=ax^2$ ではどうでしょう。\R{落}{お}ちるボール $y=5x^2$（$x$ は\R{時間}{じ|かん}〔s〕、$y$ は\R{落}{お}ちた\R{距離}{きょ|り}〔m〕）で\R{調}{しら}べます。
\begin{bkcenter}
\begin{bktable}{colspec={Q[l,m,wd=38mm] *{4}{Q[c,m,wd=15mm]}}, rowsep=1.5pt}
$x$（s） & 0 & 1 & 2 & 3 \\
$y=5x^2$（m） & 0 & 5 & 20 & 45 \\
\end{bktable}
\end{bkcenter}
$x$ が1ずつふえるたびに、$y$ の\R{変化量}{へん|か|りょう}は $5-0=5$、$20-5=15$、$45-20=25$ です。$x$ の\R{変化量}{へん|か|りょう}はどれも1なので、\R{変化}{へん|か}の\R{割合}{わり|あい}は5、15、25 で、どんどん\R{大}{おお}きくなります。$x$ が0から3までの\R{全体}{ぜん|たい}では、
\bkthrjline{（$y$ の\R{変化量}{へん|か|りょう}）÷（$x$ の\R{変化量}{へん|か|りょう}）＝$(45-0)\div(3-0)=\dfrac{45}{3}=15$}
です。\R{変化}{へん|か}の\R{割合}{わり|あい}は、どこからどこまでをくらべるかで\R{変}{か}わり、\R{一定}{いっ|てい}ではありません。

\noindent\begin{minipage}[c]{0.5\linewidth}\centering
\begin{tikzpicture}
\begin{axis}[bkplot, width=58mm, height=56mm, xmin=0, xmax=3.6, ymin=0, ymax=50, xtick={1,2,3}, ytick={5,20,45},
  xlabel={$x$（s）}, ylabel={$y$（m）}]
  \addplot[bkcurve, domain=0:3.2] {5*x^2};
  \addplot[bkOrange, line width=1pt] coordinates {(0,0) (1,5)};
  \addplot[bkOrange, line width=1pt] coordinates {(1,5) (2,20)};
  \addplot[bkOrange, line width=1pt] coordinates {(2,20) (3,45)};
  \addplot[only marks, mark=*, mark size=1.6pt, bkInk] coordinates {(0,0) (1,5) (2,20) (3,45)};
  \node[font=\sffamily\scriptsize, text=bkOrange!80!black, anchor=south east] at (axis cs:0.45,4) {5};
  \node[font=\sffamily\scriptsize, text=bkOrange!80!black, anchor=south east] at (axis cs:1.4,15) {15};
  \node[font=\sffamily\scriptsize, text=bkOrange!80!black, anchor=south east] at (axis cs:2.4,34) {25};
\end{axis}
\end{tikzpicture}
\bkfigcaption{オレンジの\R{線}{せん}の\R{傾}{かたむ}きが\R{変化}{へん|か}の\R{割合}{わり|あい}。}
\end{minipage}\hfill
\begin{minipage}[c]{0.46\linewidth}
グラフで\R{見}{み}ると、\R{変化}{へん|か}の\R{割合}{わり|あい}は、2\R{点}{てん}を\R{結}{むす}ぶ\R{直線}{ちょく|せん}（オレンジ）の\R{傾}{かたむ}きです（3-5）。\R{右}{みぎ}へ\R{行}{い}くほど、オレンジの\R{線}{せん}は\R{急}{きゅう}になります。\R{放物線}{ほう|ぶつ|せん}がだんだん\R{急}{きゅう}になっていくのは、このためです。

$x$ が\R{時間}{じ|かん}、$y$ が\R{位置}{い|ち}なので、これは $x$-$t$ グラフ（3-4）と\R{同}{おな}じ\R{形}{かたち}です。オレンジの\R{線}{せん}の\R{傾}{かたむ}きは、その1\R{秒間}{びょう|かん}の\R{平均}{へい|きん}の\R{速}{はや}さ（\qty{5}{m/s}、\qty{15}{m/s}、\qty{25}{m/s}）を\R{表}{あらわ}します（3-5）。
\end{minipage}

\bkthrjhead{あいだを\R{短}{みじか}くしていく}
こんどは、$x$ が1から $1+h$ まで\R{変}{か}わるときの\R{変化}{へん|か}の\R{割合}{わり|あい}を、\R{文字}{も|じ}のまま\R{計算}{けい|さん}してみます。$h$ は「1からどれだけ\R{進}{すす}んだか」を\R{表}{あらわ}す\R{数}{かず}で、1、0.5、0.1 など、いろいろな\R{数}{かず}が\R{入}{はい}ります。
\bkthrjstep{1}{$x=1+h$ のときの $y$ を\R{求}{もと}める}
$y=5x^2$ の $x$ に、$1+h$ をかっこごと\R{入}{い}れます（2-13）。$(1+h)^2$ は、\R{展開}{てん|かい}すると $1+2h+h^2$ です（3-8：$(a+b)^2=a^2+2ab+b^2$ の $a$ が1、$b$ が $h$）。
\[ 5(1+h)^2=5(1+2h+h^2)=5+10h+5h^2 \]
\bkthrjstep{2}{$y$ の\R{変化量}{へん|か|りょう}を\R{求}{もと}める}
$x=1$ のときの $y$ は $5\times1^2=5$ です。あと－まえで、
\[ (5+10h+5h^2)-5=10h+5h^2 \]
\bkthrjstep{3}{$x$ の\R{変化量}{へん|か|りょう} $h$ でわる}
$x$ の\R{変化量}{へん|か|りょう}は $(1+h)-1=h$ です。1つずつ $h$ でわります（2-12）。
\[ \frac{10h+5h^2}{h}=\frac{10h}{h}+\frac{5h^2}{h}=10+5h \]

\noindent \R{確}{たし}かめてみましょう。$h=1$（1から2まで）なら $10+5\times1=15$ で、\R{上}{うえ}で\R{求}{もと}めた15と\R{合}{あ}います。$h$ を\R{小}{ちい}さくしていくと、\R{次}{つぎ}のようになります。
\begin{bkcenter}
\begin{bktable}{colspec={Q[l,m,wd=36mm] *{4}{Q[c,m,wd=15mm]}}, rowsep=1.5pt}
$h$ & 1 & 0.5 & 0.1 & 0.01 \\
\R{変化}{へん|か}の\R{割合}{わり|あい} $10+5h$ & 15 & 12.5 & 10.5 & 10.05 \\
\end{bktable}
\bktabcaption{$h$ を\R{小}{ちい}さくすると、\R{変化}{へん|か}の\R{割合}{わり|あい}は10にかぎりなく\R{近}{ちか}づく。}
\end{bkcenter}
あいだを\R{短}{みじか}くするほど、\R{変化}{へん|か}の\R{割合}{わり|あい}は10に\R{近}{ちか}づいていきます。これは、「ちょうど1\R{秒}{びょう}のときの、その\R{瞬間}{しゅん|かん}の\R{速}{はや}さ」が \qty{10}{m/s} であることを\R{表}{あらわ}しています。\R{第}{だい}5\R{巻}{かん}の5-1・5-2で、この\R{考}{かんが}えを\R{使}{つか}って「\R{瞬間}{しゅん|かん}の\R{速}{はや}さ」を\R{求}{もと}め、それが\R{微分}{び|ぶん}になります。

\begin{motto}[2]
\textbf{1\R{秒}{びょう}ごとに\R{落}{お}ちる\R{距離}{きょ|り}は、5、15、25、35、…}\par
\R{落}{お}ちるボールが、はじめの1\R{秒}{びょう}、\R{次}{つぎ}の1\R{秒}{びょう}、その\R{次}{つぎ}の1\R{秒}{びょう}…に\R{落}{お}ちる\R{距離}{きょ|り}は、5、15、25、35 m で、5×1、5×3、5×5、5×7 です。1、3、5、7 と、\R{奇数}{き|すう}がならびます。400\R{年}{ねん}ほど\R{前}{まえ}に、イタリアのガリレオが、\R{坂}{さか}を\R{転}{ころ}がる\R{球}{たま}で\R{見}{み}つけたきまりです。

なぜ\R{奇数}{き|すう}になるかは、\R{展開}{てん|かい}で\R{説明}{せつ|めい}できます（3-8）。$n$ \R{秒}{びょう}から $n+1$ \R{秒}{びょう}までに\R{落}{お}ちる\R{距離}{きょ|り}は
\[ 5(n+1)^2-5n^2=5(n^2+2n+1)-5n^2=5(2n+1) \]
です。$2n+1$ は、$n=0$、1、2、3 で 1、3、5、7 と、いつも\R{奇数}{き|すう}です。
\end{motto}


\bksubsection{5}{\R{放物線}{ほう|ぶつ|せん}を\R{動}{うご}かす：$y=a(x-p)^2+q$}

\leavevmode\label{bk:3-9:ugokasu}\R{投}{な}げ\R{上}{あ}げたボールのグラフは、\R{頂点}{ちょう|てん}が\R{原点}{げん|てん}にはありません。\R{頂点}{ちょう|てん}がほかの\R{所}{ところ}にある\R{放物線}{ほう|ぶつ|せん}は、どんな\R{式}{しき}になるのでしょう。$y=x^2$ のグラフを、\R{上下}{じょう|げ}や\R{左右}{さ|ゆう}に\R{動}{うご}かして\R{考}{かんが}えます。

\bkthrjhead{\R{上}{うえ}に\R{動}{うご}かす：$y=x^2+2$}
$y=x^2$ のどの $y$ にも2をたすと、$y=x^2+2$ です。
\begin{bkcenter}
\begin{bktable}{colspec={Q[l,m,wd=24mm] *{5}{Q[c,m,wd=12mm]}}, rowsep=1.5pt}
$x$ & $-2$ & $-1$ & 0 & 1 & 2 \\
$y=x^2$ & 4 & 1 & 0 & 1 & 4 \\
$y=x^2+2$ & 6 & 3 & 2 & 3 & 6 \\
\end{bktable}
\end{bkcenter}
どの\R{点}{てん}も、ちょうど2だけ\R{上}{うえ}にきます。だから、グラフは $y=x^2$ を\R{上}{うえ}へ2\R{動}{うご}かした\R{放物線}{ほう|ぶつ|せん}で、\R{頂点}{ちょう|てん}は $(0,\ 2)$ です。$y=x^2-3$ なら、\R{下}{した}へ3\R{動}{うご}かした\R{形}{かたち}です。

\bkthrjhead{\R{右}{みぎ}に\R{動}{うご}かす：$y=(x-2)^2$}
こんどは、$x$ から2をひいてから2\R{乗}{じょう}します。\R{表}{ひょう}を\R{作}{つく}ってみましょう。
\begin{bkcenter}
\begin{bktable}{colspec={Q[l,m,wd=24mm] *{7}{Q[c,m,wd=10.5mm]}}, rowsep=1.5pt}
$x$ & $-1$ & 0 & 1 & 2 & 3 & 4 & 5 \\
$x-2$ & $-3$ & $-2$ & $-1$ & 0 & 1 & 2 & 3 \\
$y=(x-2)^2$ & 9 & 4 & 1 & \SetCell{bg=bkYellow!30} 0 & 1 & 4 & 9 \\
\end{bktable}
\bktabcaption{9、4、1、0、1、4、9 の\R{並}{なら}びが、$x=2$ を\R{中心}{ちゅう|しん}にしている。}
\end{bkcenter}
$y$ の\R{数}{かず}のならびは、$y=x^2$ の\R{表}{ひょう}（\pageref{bk:3-9:graph}ページ）と\R{同}{おな}じ 9、4、1、0、1、4、9 です。ただし、0になるのは $x=0$ ではなく、$x-2=0$ になる $x=2$ のときです。$x=3$ のときの $y$ は、$y=x^2$ の $x=1$ のときと\R{同}{おな}じ1です。どの $y$ の\R{値}{あたい}も、$y=x^2$ より\R{右}{みぎ}へ2だけずれた\R{所}{ところ}に\R{出}{で}てきます。

\noindent\begin{minipage}[c]{0.5\linewidth}\centering
\begin{tikzpicture}
\begin{axis}[bkplot, width=58mm, height=50mm, xmin=-2.6, xmax=4.8, ymin=-0.6, ymax=5.4, xtick={-2,2,4}, ytick={1,...,5},
  xlabel={$x$}, ylabel={$y$}]
  \addplot[bkthin, line width=1.2pt, domain=-2.1:2.1] {x^2};
  \addplot[bkcurve, domain=0:4] {(x-2)^2};
  \addplot[bkcurve2, domain=0.1:3.9] {(x-2)^2+1};
  \draw[-{Stealth[length=4pt]}, bkInk, line width=0.8pt, shorten <=1.5mm, shorten >=1.5mm] (axis cs:0,0) -- (axis cs:2,0);
  \draw[-{Stealth[length=4pt]}, bkInk, line width=0.8pt, shorten <=1.5mm, shorten >=1mm] (axis cs:2,0) -- (axis cs:2,1);
  \addplot[only marks, mark=*, mark size=1.7pt, bkInk] coordinates {(0,0) (2,0) (2,1)};
\end{axis}
\end{tikzpicture}
\bkfigcaption{\R{灰色}{はい|いろ} $y=x^2$、\R{青}{あお} $y=(x-2)^2$、オレンジ $y=(x-2)^2+1$。}
\end{minipage}\hfill
\begin{minipage}[c]{0.46\linewidth}
だから、$y=(x-2)^2$ のグラフは、$y=x^2$（\R{灰色}{はい|いろ}）を\R{右}{みぎ}へ2\R{動}{うご}かした\R{放物線}{ほう|ぶつ|せん}（\R{青}{あお}）です。\R{頂点}{ちょう|てん}は $(2,\ 0)$、\R{軸}{じく}は $x=2$ の\R{直線}{ちょく|せん}です。

\textbf{「ひく2」なのに\R{右}{みぎ}へ？}\quad $x-2$ が0になるのは、$x$ が2のときです。もとの\R{頂点}{ちょう|てん}の\R{役}{やく}わりを、$x=2$ がしているのです。$y=(x+3)^2$ なら、$x+3=0$ になる $x=-3$ が\R{頂点}{ちょう|てん}なので、\R{左}{ひだり}へ3\R{動}{うご}かした\R{形}{かたち}です。
\end{minipage}

\medskip
\bkthrjhead{\R{右}{みぎ}へも\R{上}{うえ}へも\R{動}{うご}かす}
$y=(x-2)^2+1$ は、$y=(x-2)^2$ を、さらに\R{上}{うえ}へ1\R{動}{うご}かした\R{放物線}{ほう|ぶつ|せん}（オレンジ）です。\R{頂点}{ちょう|てん}は $(2,\ 1)$ です。

$y=ax^2$ のときも\R{同}{おな}じです。まとめると、\R{次}{つぎ}のようになります。
\begin{taisetsu}
$y=a(x-p)^2+q$ のグラフは、$y=ax^2$ のグラフを\par
\hspace*{1\zw}\R{右}{みぎ}へ $p$、\R{上}{うえ}へ $q$ \R{動}{うご}かした\R{放物線}{ほう|ぶつ|せん}。\R{頂点}{ちょう|てん}は $(p,\ q)$、\R{軸}{じく}は\R{直線}{ちょく|せん} $x=p$。\par\vspace{2pt}
$p$ がマイナスなら\R{左}{ひだり}へ、$q$ がマイナスなら\R{下}{した}へ\R{動}{うご}かす。
\end{taisetsu}
たとえば $y=-2(x+1)^2+4$ は、$x+1=x-(-1)$ なので $p=-1$、$q=4$ です。$y=-2x^2$ を\R{左}{ひだり}へ1、\R{上}{うえ}へ4\R{動}{うご}かした、\R{下}{した}に\R{開}{ひら}いた\R{放物線}{ほう|ぶつ|せん}で、\R{頂点}{ちょう|てん} $(-1,\ 4)$ がいちばん\R{高}{たか}い\R{点}{てん}です。

\bkthrjhead{2\R{次関数}{じ|かん|すう}}
$y=(x-2)^2+1$ のかっこをはずすと（3-8）、
\[ y=(x^2-4x+4)+1=x^2-4x+5 \]
です。$x^2-4x+5$ のように、いちばん\R{高}{たか}い\R{次数}{じ|すう}が2の\R{式}{しき}を2\R{次式}{じ|しき}といいます（3-8）。$y$ が $x$ の2\R{次式}{じ|しき}で\R{表}{あらわ}される\R{関数}{かん|すう}を、\textbf{2\R{次関数}{じ|かん|すう}}といいます。\bkidx{にじかんすう}{2\R{次関数}{じ|かん|すう}}2\R{次関数}{じ|かん|すう}のグラフは、いつも\R{放物線}{ほう|ぶつ|せん}です。

では、\R{反対}{はん|たい}に $y=x^2-4x+5$ を\R{見}{み}せられたら、\R{頂点}{ちょう|てん}はどうやって\R{見}{み}つければよいでしょう。それが\R{次}{つぎ}の\R{平方完成}{へい|ほう|かん|せい}です。

\bksubsection{6}{\R{平方完成}{へい|ほう|かん|せい}で、\R{頂点}{ちょう|てん}を\R{見}{み}つける}

\leavevmode\label{bk:3-9:heihou}3-8 で、$x^2+6x+10=(x+3)^2+1$ のように、2\R{次式}{じ|しき}を「（かっこ）の2\R{乗}{じょう}＋\R{数}{かず}」の\R{形}{かたち}に\R{直}{なお}しました。これを\textbf{\R{平方完成}{へい|ほう|かん|せい}}といいました。この\R{形}{かたち}は、まさに $y=a(x-p)^2+q$ の\R{形}{かたち}です。

\bkthrjhead{\R{例}{れい}　$y=x^2+6x+10$ の\R{頂点}{ちょう|てん}}
3-8 のやり\R{方}{かた}を、もう\R{一度}{いち|ど}1\R{行}{ぎょう}ずつ\R{書}{か}きます。
\bkthrjstep{1}{$x$ の\R{係数}{けい|すう}の\R{半分}{はん|ぶん}を\R{見}{み}つける}
$x$ の\R{係数}{けい|すう}は6で、その\R{半分}{はん|ぶん}は3です。$(x+3)^2$ を\R{展開}{てん|かい}すると $x^2+6x+9$ になり、$x^2+6x$ の\R{部分}{ぶ|ぶん}がそろうからです。
\bkthrjstep{2}{9をたして、9をひく}
$+10$ の\R{前}{まえ}に、9をたして9をひきます。
\[ x^2+6x+10=x^2+6x+9-9+10=(x+3)^2+1 \]
9をたしてすぐ9をひいたので、\R{式}{しき}の\R{値}{あたい}は\R{変}{か}わっていません。
\bkthrjstep{3}{\R{頂点}{ちょう|てん}を\R{読}{よ}む}
$y=(x+3)^2+1$ は、$(x-(-3))^2+1$ なので、\R{頂点}{ちょう|てん}は $(-3,\ 1)$ です。

\bkthrjhead{なぜ、そこが\R{頂点}{ちょう|てん}（いちばん\R{低}{ひく}い\R{点}{てん}）なの？}
\noindent\begin{minipage}[c]{0.52\linewidth}
\R{理由}{り|ゆう}は2つの\R{方法}{ほう|ほう}で\R{確}{たし}かめられます。

\textbf{その1：2\R{乗}{じょう}は0より\R{小}{ちい}さくならない。}\par\noindent $(x+3)^2$ は、どんな $x$ でも0\R{以上}{い|じょう}です。0になるのは、$x+3=0$ の $x=-3$ のときだけです。だから $y=(x+3)^2+1$ は、いつも1\R{以上}{い|じょう}で、ちょうど1になるのは $x=-3$ のときです。

\textbf{その2：\R{表}{ひょう}で\R{確}{たし}かめる。}\par\noindent もとの\R{式}{しき} $y=x^2+6x+10$ に\R{数}{かず}を\R{入}{い}れると、\R{右}{みぎ}の\R{表}{ひょう}のように、$x=-3$ の1がいちばん\R{小}{ちい}さく、\R{左右}{さ|ゆう}に\R{同}{おな}じ\R{数}{かず}がならびます。
\end{minipage}\hfill
\begin{minipage}[c]{0.44\linewidth}\centering
\begin{bktable}{colspec={Q[c,m,wd=10mm] Q[c,m,wd=34mm]}, rowsep=1pt, cells={font=\small}}
$x$ & $x^2+6x+10$ \\
$-5$ & $25-30+10=5$ \\
$-4$ & $16-24+10=2$ \\
\SetCell{bg=bkYellow!30} $-3$ & \SetCell{bg=bkYellow!30} $9-18+10=1$ \\
$-2$ & $4-12+10=2$ \\
$-1$ & $1-6+10=5$ \\
\end{bktable}
\end{minipage}

\noindent\begin{minipage}[c]{0.5\linewidth}\centering
\begin{tikzpicture}
\begin{axis}[bkplot, width=62mm, height=62mm, xmin=-6.4, xmax=1.4, ymin=-0.6, ymax=6.4, xtick={-5,-4,-3,-2,-1}, ytick={1,2,5},
  xlabel={$x$}, ylabel={$y$}]
  \addplot[bkcurve, domain=-5.2:-0.8] {(x+3)^2+1};
  \addplot[bkRed, line width=0.9pt, dash pattern=on 3pt off 2pt] coordinates {(-3,0) (-3,6)};
  \addplot[only marks, mark=*, mark size=1.5pt, bkInk] coordinates {(-5,5) (-4,2) (-2,2) (-1,5)};
  \addplot[only marks, mark=*, mark size=2.1pt, bkRed] coordinates {(-3,1)};
\end{axis}
\end{tikzpicture}
\bkfigcaption{$y=x^2+6x+10=(x+3)^2+1$}
\end{minipage}\hfill
\begin{minipage}[c]{0.46\linewidth}
グラフにすると、\R{頂点}{ちょう|てん} $(-3,\ 1)$（\R{赤}{あか}）がいちばん\R{低}{ひく}い\R{点}{てん}で、\R{軸}{じく}は\R{直線}{ちょく|せん} $x=-3$（\R{赤}{あか}い\R{点線}{てん|せん}）です。$y=x^2$ を\R{左}{ひだり}へ3、\R{上}{うえ}へ1\R{動}{うご}かした\R{放物線}{ほう|ぶつ|せん}です。\R{式}{しき}を\R{見}{み}ただけでは\R{頂点}{ちょう|てん}がわからなくても、\R{平方完成}{へい|ほう|かん|せい}すれば、\R{計算}{けい|さん}でわかるのです。
\end{minipage}

\bkthrjhead{$x^2$ の\R{係数}{けい|すう}が1でないとき：\R{投}{な}げ\R{上}{あ}げたボール}
\noindent\begin{minipage}[c]{0.6\linewidth}
ボールを\R{真上}{ま|うえ}に、\qty{20}{m/s} の\R{速}{はや}さで\R{投}{な}げ\R{上}{あ}げます（\R{空気}{くう|き}の\R{抵抗}{てい|こう}は\R{考}{かんが}えません）。もし\R{重力}{じゅう|りょく}がなければ、ボールは\R{同}{おな}じ\R{速}{はや}さで\R{上}{あ}がりつづけ、$t$ \R{秒後}{びょう|ご}には $20t$ m の\R{高}{たか}さにいるはずです（\R{道}{みち}のり＝\R{速}{はや}さ×\R{時間}{じ|かん}）。でも、\R{重力}{じゅう|りょく}があるので、ボールはその\R{高}{たか}さから、\R{落}{お}ちる\R{距離}{きょ|り}の $5t^2$ m だけ\R{下}{した}にきます（2-13）。だから、$t$ \R{秒後}{びょう|ご}の\R{高}{たか}さ $y$〔m〕は、
\[ y=20t-5t^2=-5t^2+20t \]
です。$t$ に0から4までを\R{入}{い}れると、\R{高}{たか}さは\R{右}{みぎ}の\R{表}{ひょう}のようになります。2\R{秒}{びょう}のときの20 m がいちばん\R{高}{たか}そうです。
\end{minipage}\hfill
\begin{minipage}[c]{0.36\linewidth}\centering
\begin{bktable}{colspec={Q[c,m,wd=10mm] Q[c,m,wd=24mm]}, rowsep=1pt, cells={font=\small}}
$t$（s） & $20t-5t^2$（m） \\
0 & $0-0=0$ \\
1 & $20-5=15$ \\
2 & $40-20=20$ \\
3 & $60-45=15$ \\
4 & $80-80=0$ \\
\end{bktable}
\end{minipage}

\medskip
\R{平方完成}{へい|ほう|かん|せい}で\R{確}{たし}かめましょう。$t^2$ の\R{前}{まえ}に $-5$ がついているので、1つ\R{手順}{て|じゅん}がふえます。
\bkthrjstep{1}{$t^2$ の\R{係数}{けい|すう} $-5$ でくくる}
$-5t^2+20t$ の2つの\R{項}{こう}を、どちらも $-5$ でわった\R{形}{かたち}にします（3-8：\R{共通因数}{きょう|つう|いん|すう}でくくる）。
\[ -5t^2+20t=-5(t^2-4t) \]
\R{確}{たし}かめ：かっこをはずすと $-5\times t^2=-5t^2$、$-5\times(-4t)=20t$ で、もとにもどります。
\bkthrjstep{2}{かっこの\R{中}{なか}を\R{平方完成}{へい|ほう|かん|せい}する}
$t$ の\R{係数}{けい|すう} $-4$ の\R{半分}{はん|ぶん}は $-2$ です。$(t-2)^2=t^2-4t+4$ なので、4をたして4をひきます。
\[ t^2-4t=t^2-4t+4-4=(t-2)^2-4 \]
\bkthrjstep{3}{もとの\R{式}{しき}にもどす}
②の\R{結果}{けっ|か}を、①のかっこの\R{中}{なか}に\R{入}{い}れます。
\[ y=-5\{(t-2)^2-4\}=-5(t-2)^2+(-5)\times(-4)=-5(t-2)^2+20 \]
$-5$ を、$(t-2)^2$ にも $-4$ にもかけるのを\R{忘}{わす}れないようにします（\R{分配法則}{ぶん|ぱい|ほう|そく}、2-12）。

\noindent\begin{minipage}[c]{0.5\linewidth}\centering
\begin{tikzpicture}
\begin{axis}[bkplot, width=58mm, height=50mm, xmin=0, xmax=4.8, ymin=0, ymax=25, xtick={1,2,3,4}, ytick={5,10,15,20},
  xlabel={$t$（s）}, ylabel={$y$（m）}]
  \addplot[bkcurve, domain=0:4] {20*x-5*x^2};
  \addplot[bkRed, line width=0.9pt, dash pattern=on 3pt off 2pt] coordinates {(2,0) (2,23.5)};
  \addplot[only marks, mark=*, mark size=1.6pt, bkInk] coordinates {(0,0) (1,15) (3,15) (4,0)};
  \addplot[only marks, mark=*, mark size=2.1pt, bkRed] coordinates {(2,20)};
  \node[font=\sffamily\scriptsize, text=bkRed!80!black, anchor=west] at (axis cs:2.15,22.6) {\R{頂点}{ちょう|てん} $(2,\ 20)$};
\end{axis}
\end{tikzpicture}
\bkfigcaption{$y=-5(t-2)^2+20$ のグラフ。}
\end{minipage}\hfill
\begin{minipage}[c]{0.46\linewidth}
$y=-5(t-2)^2+20$ は、\R{頂点}{ちょう|てん}が $(2,\ 20)$ の、\R{下}{した}に\R{開}{ひら}いた\R{放物線}{ほう|ぶつ|せん}です。$(t-2)^2$ は0\R{以上}{い|じょう}なので、それに $-5$ をかけた $-5(t-2)^2$ は0\R{以下}{い|か}です。だから $y$ は20をこえず、ちょうど20になるのは $t=2$ のときです。

ボールは、\R{投}{な}げてから2\R{秒後}{びょう|ご}に、いちばん\R{高}{たか}い \qty{20}{m} に\R{届}{とど}きます。\R{表}{ひょう}で\R{見}{み}た\R{答}{こた}えと\R{合}{あ}いました。
\end{minipage}

\medskip
\bkthrjhead{\R{注意}{ちゅう|い}：このグラフは、ボールの\R{通}{とお}り\R{道}{みち}ではない}
\R{真上}{ま|うえ}に\R{投}{な}げたボールは、まっすぐ\R{上}{あ}がって、まっすぐ\R{下}{お}りてきます。\R{上}{うえ}のグラフが\R{山}{やま}の\R{形}{かたち}なのは、\R{横軸}{よこ|じく}が「\R{時刻}{じ|こく}」だからです。ボールが\R{横}{よこ}に\R{動}{うご}いているわけではありません。

\bksubsection{7}{\R{範囲}{はん|い}があるときの、\R{最大}{さい|だい}と\R{最小}{さい|しょう}}

\leavevmode\label{bk:3-9:hani}$x$ の\R{範囲}{はん|い}（\R{変域}{へん|いき}、3-1）が\R{決}{き}まっているときは、グラフのうち、その\R{範囲}{はん|い}の\R{部分}{ぶ|ぶん}だけを\R{見}{み}ます。いちばん\R{大}{おお}きい $y$ の\R{値}{あたい}を\textbf{\R{最大値}{さい|だい|ち}}、いちばん\R{小}{ちい}さい\R{値}{あたい}を\textbf{\R{最小値}{さい|しょう|ち}}といいます。\bkidx{さいだいち・さいしょうち}{\R{最大値}{さい|だい|ち}・\R{最小値}{さい|しょう|ち}}

\noindent\begin{minipage}[t]{0.48\linewidth}\centering
\begin{tikzpicture}
\begin{axis}[bkplot, width=54mm, height=54mm, xmin=-2.4, xmax=3.6, ymin=-0.8, ymax=10.4, xtick={-2,-1,1,2,3}, ytick={1,4,9},
  xlabel={$x$}, ylabel={$y$}]
  \addplot[bkthin, line width=1.1pt, domain=-2.2:3.2] {x^2};
  \addplot[bkcurve, line width=2.4pt, domain=1:3] {x^2};
  \addplot[only marks, mark=*, mark size=1.9pt, bkInk] coordinates {(1,1) (3,9)};
\end{axis}
\end{tikzpicture}
\bkfigcaption{$1\leqq x\leqq3$ のとき。\R{最大}{さい|だい}9、\R{最小}{さい|しょう}1。}
\end{minipage}\hfill
\begin{minipage}[t]{0.48\linewidth}\centering
\begin{tikzpicture}
\begin{axis}[bkplot, width=54mm, height=54mm, xmin=-2.4, xmax=3.6, ymin=-0.8, ymax=10.4, xtick={-2,-1,1,2,3}, ytick={1,4,9},
  xlabel={$x$}, ylabel={$y$}]
  \addplot[bkthin, line width=1.1pt, domain=-2.2:3.2] {x^2};
  \addplot[bkcurve, line width=2.4pt, domain=-1:2] {x^2};
  \addplot[only marks, mark=*, mark size=1.9pt, bkInk] coordinates {(-1,1) (0,0) (2,4)};
\end{axis}
\end{tikzpicture}
\bkfigcaption{$-1\leqq x\leqq2$ のとき。\R{最大}{さい|だい}4、\R{最小}{さい|しょう}0。}
\end{minipage}

\medskip
\textbf{$y=x^2$ で、$1\leqq x\leqq3$ のとき}（\R{左}{ひだり}の\R{図}{ず}）\par
この\R{範囲}{はん|い}には、\R{頂点}{ちょう|てん} $x=0$ が\R{入}{はい}っていません。$x$ が1から3へふえると、$y$ は1から9へ、ずっとふえます。だから、\R{最小値}{さい|しょう|ち}は $x=1$ のときの1、\R{最大値}{さい|だい|ち}は $x=3$ のときの9です。

\textbf{$y=x^2$ で、$-1\leqq x\leqq2$ のとき}（\R{右}{みぎ}の\R{図}{ず}）\par
こんどは、\R{頂点}{ちょう|てん} $x=0$ が\R{範囲}{はん|い}の\R{中}{なか}にあります。\R{下}{さ}がってから\R{上}{あ}がるので、いちばん\R{低}{ひく}いのは\R{頂点}{ちょう|てん}で、\R{最小値}{さい|しょう|ち}は0です。いちばん\R{高}{たか}いのは、はしのどちらかです。$x=-1$ では1、$x=2$ では4なので、\R{最大値}{さい|だい|ち}は $x=2$ のときの4です。$x=-1$ のときの1を\R{最小値}{さい|しょう|ち}だと\R{思}{おも}ってしまわないように\R{注意}{ちゅう|い}しましょう。

\begin{taisetsu}
\R{範囲}{はん|い}があるときの\R{最大}{さい|だい}・\R{最小}{さい|しょう}は、\R{次}{つぎ}の\R{順}{じゅん}にさがす。\par
①\enspace \R{頂点}{ちょう|てん}が\R{範囲}{はん|い}の\R{中}{なか}にあるかを\R{見}{み}る。\par
②\enspace \R{範囲}{はん|い}のはしの2つの\R{値}{あたい}と、（\R{中}{なか}にあれば）\R{頂点}{ちょう|てん}の\R{値}{あたい}を\R{計算}{けい|さん}する。\par
③\enspace その\R{中}{なか}で、いちばん\R{大}{おお}きいのが\R{最大値}{さい|だい|ち}、いちばん\R{小}{ちい}さいのが\R{最小値}{さい|しょう|ち}。
\end{taisetsu}

\bksubsection{8}{はてな？の\R{答}{こた}え：\R{横}{よこ}は\R{比例}{ひ|れい}、たては2\R{乗}{じょう}に\R{比例}{ひ|れい}}

\leavevmode\label{bk:3-9:kotae}\bkhatenaref のボールにもどりましょう。やってみようの\R{表}{ひょう}から、2つのことがわかりました。
\begin{itemize}[leftmargin=1.6\zw, itemsep=1pt, topsep=2pt, label={\color{\bkTT}●}]
\item \R{横}{よこ}は、\qty{2}{m/s} の\R{一定}{いっ|てい}の\R{速}{はや}さで\R{進}{すす}む：$x=2t$\par\noindent （\R{飛}{と}び\R{出}{だ}したあと、\R{横}{よこ}にはおす\R{力}{ちから}も\R{引}{ひ}く\R{力}{ちから}もはたらかないので、\R{横}{よこ}の\R{速}{はや}さは\R{変}{か}わりません。\R{空気}{くう|き}の\R{抵抗}{てい|こう}は\R{考}{かんが}えません。）
\item たては、\R{落}{お}ちるボールと\R{同}{おな}じに\R{下}{さ}がる：$y=5t^2$
\end{itemize}
\begin{bkfigure}
\begin{tikzpicture}[x=1mm, y=1mm]
  % 1 m を 60 mm。横 0.2 m ごと（12 mm）、下がる 0, 0.05, 0.2, 0.45, 0.8 m（0, 3, 12, 27, 48 mm）
  \foreach \k/\d in {0/0, 1/3, 2/12, 3/27, 4/48} {
    \draw[bkGray, line width=0.4pt, densely dotted] (12*\k,6) -- (12*\k,-\d);
    \draw[bkGray, line width=0.4pt, densely dotted] (-6,-\d) -- (12*\k,-\d);
    \fill[bkBlue] (12*\k,6) circle (1.2);
    \fill[bkRed] (-6,-\d) circle (1.2);
    \bkthrjball{(12*\k,-\d)}{1.8}}
  \draw[bkBlue, line width=0.8pt] (-2,6) -- (54,6);
  \draw[bkRed, line width=0.8pt] (-6,2) -- (-6,-52);
  \foreach \k/\lab in {1/0.2, 2/0.4, 3/0.6, 4/0.8} {\node[font=\sffamily\scriptsize, text=bkBlue!80!black, anchor=south] at (12*\k,8) {\lab};}
  \foreach \d/\lab in {3/0.05, 12/0.2, 27/0.45, 48/0.8} {\node[font=\sffamily\scriptsize, text=bkRed!80!black, anchor=east] at (-8.5,-\d) {\lab};}
  \node[font=\sffamily\footnotesize, text=bkBlue!80!black, anchor=west, align=left] at (60,6) {\R{横}{よこ}のかげ：\\\R{同}{おな}じ\R{間隔}{かん|かく}でならぶ\\（\R{比例}{ひ|れい}）};
  \node[font=\sffamily\footnotesize, text=bkRed!80!black, anchor=west, align=left] at (60,-30) {たてのかげ：\\\R{間隔}{かん|かく}が1、3、5、7\\の\R{割合}{わり|あい}で\R{広}{ひろ}がる\\（2\R{乗}{じょう}に\R{比例}{ひ|れい}）};
\end{tikzpicture}
\bkfigcaption{ボールの\R{位置}{い|ち}を、\R{上}{うえ}（\R{青}{あお}）と\R{左}{ひだり}（\R{赤}{あか}）にうつした「かげ」。\R{数}{かず}の\R{単位}{たん|い}は m。}
\end{bkfigure}
ボールの\R{位置}{い|ち}を、\R{真上}{ま|うえ}と\R{真横}{ま|よこ}から\R{見}{み}た「かげ」に\R{分}{わ}けてみましょう。\R{上}{うえ}のかげ（\R{青}{あお}）は、\R{同}{おな}じ\R{間隔}{かん|かく}でならびます。\R{左}{ひだり}のかげ（\R{赤}{あか}）は、0.05、0.2、0.45、0.8 と、\R{間隔}{かん|かく}がだんだん\R{広}{ひろ}がります。2つのかげを\R{合}{あ}わせた\R{点}{てん}が、ボールの\R{位置}{い|ち}です。

この2つの\R{式}{しき}から $t$ を\R{消}{け}して、$x$ と $y$ だけの\R{関係}{かん|けい}にします。

\bkthrjstep{1}{$x=2t$ から、$t$ を $x$ の\R{式}{しき}で\R{表}{あらわ}す}
\R{両辺}{りょう|へん}を2でわると（2-14）、$t=\dfrac{x}{2}$ です。「\R{横}{よこ}に $x$〔m〕\R{進}{すす}んだのは、$\dfrac{x}{2}$ \R{秒後}{びょう|ご}」という\R{意味}{い|み}です。
\bkthrjstep{2}{$y=5t^2$ の $t$ に、かっこごと\R{入}{い}れる}
$t$ のところに $\dfrac{x}{2}$ を\R{入}{い}れます（2-13）。\R{分数}{ぶん|すう}の2\R{乗}{じょう}は、\R{分子}{ぶん|し}も\R{分母}{ぶん|ぼ}も2\R{乗}{じょう}します。
\[ y=5\left(\frac{x}{2}\right)^2=5\times\frac{x^2}{4}=\frac{5}{4}x^2=1.25x^2 \]
\bkthrjstep{3}{\R{表}{ひょう}で\R{確}{たし}かめる}
\R{表}{ひょう}の $x$ を\R{入}{い}れてみます。
\begin{align*}
x=0.4\text{ のとき}\quad 1.25\times0.4^2&=1.25\times0.16=0.2\\
x=0.8\text{ のとき}\quad 1.25\times0.8^2&=1.25\times0.64=0.8
\end{align*}
どちらも、\R{表}{ひょう}の\R{下}{さ}がった\R{長}{なが}さと\R{合}{あ}います。\R{表}{ひょう}の $y\div x^2$ も、$0.05\div0.04=1.25$、$0.20\div0.16=1.25$、… と、いつも1.25 です。

\noindent\begin{minipage}[c]{0.52\linewidth}
\R{下}{さ}がった\R{長}{なが}さ $y$ は、\R{横}{よこ} $x$ の2\R{乗}{じょう}に\R{比例}{ひ|れい}しています。だから、ボールの\R{通}{とお}り\R{道}{みち}は\R{放物線}{ほう|ぶつ|せん}です。

\R{図}{ず}では\R{下向}{した|む}きを＋にしましたが、\R{上向}{うわ|む}きを＋にすると（＋の\R{向}{む}きは\R{自分}{じ|ぶん}で\R{決}{き}めてよい。3-2）、\R{通}{とお}り\R{道}{みち}は $y=-1.25x^2$ で、\R{下}{した}に\R{開}{ひら}いた\R{放物線}{ほう|ぶつ|せん}の\R{右半分}{みぎ|はん|ぶん}になります。\R{机}{つくえ}を\R{飛}{と}び\R{出}{だ}した\R{所}{ところ}が、\R{放物線}{ほう|ぶつ|せん}の\R{頂点}{ちょう|てん}です。その\R{瞬間}{しゅん|かん}、ボールは\R{真横}{ま|よこ}に\R{動}{うご}いていて、\R{上}{うえ}にも\R{下}{した}にも\R{動}{うご}いていないからです。
\end{minipage}\hfill
\begin{minipage}[c]{0.44\linewidth}\centering
\begin{tikzpicture}
\begin{axis}[bkplot, width=52mm, height=50mm, xmin=0, xmax=1.05, ymin=-0.98, ymax=0.1, xtick={0.2,0.4,0.6,0.8}, ytick={-0.8,-0.45,-0.2},
  xticklabel style={/pgf/number format/fixed}, yticklabel style={/pgf/number format/fixed},
  axis x line=top, axis y line=left, xlabel={$x$（m）}, ylabel={$y$（m）},
  every axis y label/.style={at={(ticklabel* cs:0)}, anchor=north, yshift=-1mm, font=\sffamily\small, text=bkInk}]
  \addplot[bkcurve, domain=0:0.85] {-1.25*x^2};
  \addplot[only marks, mark=*, mark size=1.7pt, bkInk] coordinates {(0,0) (0.2,-0.05) (0.4,-0.2) (0.6,-0.45) (0.8,-0.8)};
  \node[font=\sffamily\scriptsize, text=bkBlue!80!black, anchor=north east] at (axis cs:0.55,-0.5) {$y=-1.25x^2$};
\end{axis}
\end{tikzpicture}
\end{minipage}

\medskip
\bkhatenaref の\R{答}{こた}え：ボールが\R{横}{よこ}に\R{一定}{いっ|てい}の\R{速}{はや}さで\R{進}{すす}みながら、2\R{乗}{じょう}に\R{比例}{ひ|れい}して\R{落}{お}ちていくので、\R{通}{とお}り\R{道}{みち}は\R{放物線}{ほう|ぶつ|せん}になります。ななめ\R{上}{うえ}に\R{投}{な}げたときも、\R{横}{よこ}は\R{一定}{いっ|てい}の\R{速}{はや}さ、たては\R{投}{な}げ\R{上}{あ}げと\R{同}{おな}じ\R{式}{しき}になるので、\R{通}{とお}り\R{道}{みち}はやはり\R{放物線}{ほう|ぶつ|せん}です。そのときは、\R{山}{やま}の\R{頂点}{ちょう|てん}が、ボールのいちばん\R{高}{たか}い\R{所}{ところ}になります。

\begin{kotoba}
\small
\bkword{2\R{乗}{じょう}に\R{比例}{ひ|れい}}{$y=ax^2$ の\R{関係}{かん|けい}。$a=\dfrac{y}{x^2}$ が\R{比例定数}{ひ|れい|てい|すう}。$x$ が2\R{倍}{ばい}で $y$ は4\R{倍}{ばい}。}
\bkword{\R{放物線}{ほう|ぶつ|せん}}{$y=ax^2$ や2\R{次関数}{じ|かん|すう}のグラフの\R{形}{かたち}。\R{物}{もの}を\R{放}{はな}ったときの\R{通}{とお}り\R{道}{みち}の\R{形}{かたち}。}
\bkword{\R{軸}{じく}・\R{頂点}{ちょう|てん}}{\R{放物線}{ほう|ぶつ|せん}を\R{折}{お}って\R{重}{かさ}なる\R{折}{お}り\R{目}{め}の\R{直線}{ちょく|せん}が\R{軸}{じく}。\R{軸}{じく}と\R{放物線}{ほう|ぶつ|せん}の\R{交}{まじ}わる\R{点}{てん}が\R{頂点}{ちょう|てん}。}
\bkword{2\R{次関数}{じ|かん|すう}}{$y$ が $x$ の2\R{次式}{じ|しき}で\R{表}{あらわ}される\R{関数}{かん|すう}。グラフは\R{放物線}{ほう|ぶつ|せん}。}
\bkword{\R{最大値}{さい|だい|ち}・\R{最小値}{さい|しょう|ち}}{\R{範囲}{はん|い}の\R{中}{なか}で、いちばん\R{大}{おお}きい\R{値}{あたい}・いちばん\R{小}{ちい}さい\R{値}{あたい}。}
\end{kotoba}

%% ---------------- 例題 ----------------
\Needspace{14\baselineskip}
\bksection{reidai}{\R{例題}{れい|だい}でたしかめよう}

\begin{reidai}[2\R{乗}{じょう}に\R{比例}{ひ|れい}する\R{式}{しき}を\R{作}{つく}る]
$y$ は $x$ の2\R{乗}{じょう}に\R{比例}{ひ|れい}し、$x=2$ のとき $y=12$ です。$y$ を $x$ の\R{式}{しき}で\R{表}{あらわ}し、$x=5$ のときと $x=-1$ のときの $y$ を\R{求}{もと}めましょう。\par
\kangaekata 2\R{乗}{じょう}に\R{比例}{ひ|れい}するので、$y=ax^2$ の\R{形}{かたち}です。\R{比例定数}{ひ|れい|てい|すう}は $a=\dfrac{y}{x^2}$ で\R{求}{もと}めます。\par
\shiki $a=\dfrac{12}{2^2}=\dfrac{12}{4}=3$\quad わかっている $x=2$、$y=12$ を\R{使}{つか}う\par
\hspace*{2.6\zw}だから $y=3x^2$\par
\hspace*{2.6\zw}$x=5$ のとき $y=3\times5^2=3\times25=75$、$x=-1$ のとき $y=3\times(-1)^2=3$\par
\kotae $y=3x^2$。$x=5$ のとき $y=75$、$x=-1$ のとき $y=3$\par
（\R{確}{たし}かめ：$x$ が2から4へ2\R{倍}{ばい}なら、$y$ は $3\times16=48$ で、12の4\R{倍}{ばい}。）
\end{reidai}

\begin{reidai}[\R{表}{ひょう}を\R{作}{つく}って、グラフをかく]
\R{次}{つぎ}の\R{関数}{かん|すう}について、$x=-2$、$-1$、0、1、2 のときの $y$ を\R{求}{もと}めて\R{表}{ひょう}にし、グラフをノートにかきましょう。\par（1）\nobreak $y=\dfrac{1}{2}x^2$\quad（2）\nobreak $y=-\dfrac{1}{2}x^2$\par
\kangaekata $x$ の2\R{乗}{じょう}を\R{先}{さき}に\R{計算}{けい|さん}してから、$a$ をかけます。（2）は、（1）の $y$ にマイナスをつけた\R{値}{あたい}になります。\par
\shiki（1）\nobreak $x=2$ のとき\quad $\dfrac{1}{2}\times2^2=\dfrac{1}{2}\times4=2$\par\vspace{2pt}
\hspace*{4.3\zw}$x=-1$ のとき\quad $\dfrac{1}{2}\times(-1)^2=\dfrac{1}{2}\times1=0.5$\par\vspace{2pt}
\hspace*{2.6\zw}（2）\nobreak $x=2$ のとき\quad $-\dfrac{1}{2}\times4=-2$\par\vspace{2pt}
\hspace*{4.3\zw}$x=-1$ のとき\quad $-\dfrac{1}{2}\times1=-0.5$\par\vspace{2pt}
\hspace*{4.3\zw}ほかの $x$ も\R{同}{おな}じ。$x=2$ と $x=-2$ は、\R{同}{おな}じ $y$ になる。\par
\noindent\begin{minipage}[c]{0.56\linewidth}
\begin{bktable}{colspec={Q[l,m,wd=15mm] *{5}{Q[c,m,wd=8.5mm]}}, rowsep=1pt, colsep=2pt, cells={font=\small}}
$x$ & $-2$ & $-1$ & 0 & 1 & 2 \\
（1）$y$ & 2 & 0.5 & 0 & 0.5 & 2 \\
（2）$y$ & $-2$ & $-0.5$ & 0 & $-0.5$ & $-2$ \\
\end{bktable}
\end{minipage}\hfill
\begin{minipage}[c]{0.42\linewidth}\centering
\begin{tikzpicture}
\begin{axis}[bkplot, width=46mm, height=44mm, xmin=-2.6, xmax=3.2, ymin=-2.6, ymax=2.6, xtick={-2,2}, ytick={-2,2},
  xlabel={$x$}, ylabel={$y$}]
  \addplot[bkcurve, domain=-2.2:2.2] {0.5*x^2};
  \addplot[bkRed, line width=1.6pt, domain=-2.2:2.2] {-0.5*x^2};
  \node[font=\sffamily\scriptsize, text=bkBlue!80!black, anchor=west] at (axis cs:2.25,2.2) {（1）};
  \node[font=\sffamily\scriptsize, text=bkRed!80!black, anchor=west] at (axis cs:2.25,-2.2) {（2）};
\end{axis}
\end{tikzpicture}
\end{minipage}\par
\kotae \R{表}{ひょう}と\R{図}{ず}のとおり。（1）\nobreak は\R{上}{うえ}に\R{開}{ひら}いた\R{放物線}{ほう|ぶつ|せん}、（2）\nobreak は\R{下}{した}に\R{開}{ひら}いた\R{放物線}{ほう|ぶつ|せん}で、どちらも\R{頂点}{ちょう|てん}は\R{原点}{げん|てん}、\R{軸}{じく}は $y$ \R{軸}{じく}。（2）\nobreak は（1）\nobreak を $x$ \R{軸}{じく}で\R{折}{お}り\R{返}{かえ}した\R{形}{かたち}。
\end{reidai}

\begingroup\tcbset{equal height group=bkthrjA}%
\noindent\begin{minipage}[t]{0.49\linewidth}
\begin{jibun}
\raggedright $y$ は $x$ の2\R{乗}{じょう}に\R{比例}{ひ|れい}し、$x=3$ のとき $y=18$ です。$y$ を $x$ の\R{式}{しき}で\R{表}{あらわ}し、$x=4$ のときと $x=-2$ のときの $y$ を\R{求}{もと}めましょう。
\bkjans{$a=\dfrac{18}{3^2}=2$ で $y=2x^2$。$x=4$ のとき $y=2\times16=32$、$x=-2$ のとき $y=2\times4=8$}
\end{jibun}
\end{minipage}\hfill
\begin{minipage}[t]{0.49\linewidth}
\begin{jibun}
\raggedright \R{次}{つぎ}の\R{関数}{かん|すう}について、$x=-2$、$-1$、0、1、2 のときの $y$ を\R{求}{もと}めて\R{表}{ひょう}にし、グラフをノートにかきましょう。\par（1）\nobreak $y=\dfrac{1}{4}x^2$\par（2）\nobreak $y=-2x^2$
\bkjans{（1）\nobreak 1、0.25、0、0.25、1。\R{上}{うえ}に\R{開}{ひら}いた\R{放物線}{ほう|ぶつ|せん}。（2）\nobreak $-8$、$-2$、0、$-2$、$-8$。\R{下}{した}に\R{開}{ひら}いた\R{放物線}{ほう|ぶつ|せん}。どちらも\R{頂点}{ちょう|てん}は\R{原点}{げん|てん}、\R{軸}{じく}は $y$ \R{軸}{じく}}
\end{jibun}
\end{minipage}\par
\endgroup

\begin{reidai}[\R{変化}{へん|か}の\R{割合}{わり|あい}]
$y=2x^2$ について、\R{次}{つぎ}の\R{変化}{へん|か}の\R{割合}{わり|あい}を\R{求}{もと}めましょう。\par（1）\nobreak $x$ が1から3まで\quad（2）\nobreak $x$ が1から $1+h$ まで\par
\kangaekata \R{変化}{へん|か}の\R{割合}{わり|あい}＝$y$ の\R{変化量}{へん|か|りょう}÷$x$ の\R{変化量}{へん|か|りょう}（3-5）。（2）は、$1+h$ をかっこごと\R{入}{い}れ、$(1+h)^2=1+2h+h^2$（3-8）を\R{使}{つか}います。\par
\shiki（1）\nobreak $x=1$ で $y=2$、$x=3$ で $y=2\times9=18$\par
\addvspace{4pt}\hspace*{4.3\zw}$\dfrac{18-2}{3-1}=\dfrac{16}{2}=8$\par\vspace{2pt}
\hspace*{2.6\zw}（2）\nobreak $2(1+h)^2=2(1+2h+h^2)=2+4h+2h^2$\par
\hspace*{4.3\zw}$y$ の\R{変化量}{へん|か|りょう}は $(2+4h+2h^2)-2=4h+2h^2$、$x$ の\R{変化量}{へん|か|りょう}は $h$\par
\addvspace{4pt}\hspace*{4.3\zw}$\dfrac{4h+2h^2}{h}=\dfrac{4h}{h}+\dfrac{2h^2}{h}=4+2h$\par\vspace{2pt}
\hspace*{4.3\zw}\R{確}{たし}かめ：$h=2$ なら $4+2\times2=8$ で、（1）と\R{合}{あ}う。\par
\kotae（1）\nobreak 8\quad（2）\nobreak $4+2h$
\end{reidai}

\begin{reidai}[\R{頂点}{ちょう|てん}と\R{軸}{じく}を\R{読}{よ}む]
\R{次}{つぎ}の\R{放物線}{ほう|ぶつ|せん}の\R{頂点}{ちょう|てん}と\R{軸}{じく}を\R{答}{こた}え、\R{上}{うえ}に\R{開}{ひら}くか\R{下}{した}に\R{開}{ひら}くかを\R{答}{こた}えましょう。\par（1）\nobreak $y=(x-3)^2+2$\quad（2）\nobreak $y=-2(x+1)^2+4$\par
\kangaekata $y=a(x-p)^2+q$ の\R{形}{かたち}とくらべて、$p$ と $q$ を\R{読}{よ}みます。$x+1$ は $x-(-1)$ と\R{見}{み}ます。\R{開}{ひら}き\R{方}{かた}は $a$ の\R{符号}{ふ|ごう}で\R{決}{き}まります。\par
\shiki（1）\nobreak $a=1$、$p=3$、$q=2$\par
\hspace*{2.6\zw}（2）\nobreak $x+1=x-(-1)$ なので、$a=-2$、$p=-1$、$q=4$\par
\kotae（1）\nobreak \R{頂点}{ちょう|てん} $(3,\ 2)$、\R{軸}{じく}は\R{直線}{ちょく|せん} $x=3$、\R{上}{うえ}に\R{開}{ひら}く\par
\hspace*{2.6\zw}（2）\nobreak \R{頂点}{ちょう|てん} $(-1,\ 4)$、\R{軸}{じく}は\R{直線}{ちょく|せん} $x=-1$、\R{下}{した}に\R{開}{ひら}く\par
\noindent\begin{minipage}[c]{0.5\linewidth}\centering
\begin{tikzpicture}
\begin{axis}[bkplot, width=64mm, height=52mm, xmin=-4.6, xmax=5.6, ymin=-4.6, ymax=8.4, xtick={-3,-1,1,3,5}, ytick=\empty,
  xlabel={$x$}, ylabel={$y$}]
  \addplot[bkcurve, domain=0.6:5.4] {(x-3)^2+2};
  \addplot[bkRed, line width=1.6pt, domain=-2.8:0.8] {-2*(x+1)^2+4};
  \addplot[only marks, mark=*, mark size=1.8pt, bkInk] coordinates {(3,2) (-1,4)};
  \node[font=\sffamily\scriptsize, text=bkBlue!80!black, anchor=north west] at (axis cs:3.2,1.7) {（1）$(3,\ 2)$};
  \node[font=\sffamily\scriptsize, text=bkRed!80!black, anchor=south east] at (axis cs:-1.15,4.4) {（2）$(-1,\ 4)$};
\end{axis}
\end{tikzpicture}
\end{minipage}\hfill
\begin{minipage}[c]{0.46\linewidth}
（1）\nobreak は $y=x^2$ を\R{右}{みぎ}へ3、\R{上}{うえ}へ2\R{動}{うご}かした\R{形}{かたち}（\R{青}{あお}）。（2）\nobreak は $y=-2x^2$ を\R{左}{ひだり}へ1、\R{上}{うえ}へ4\R{動}{うご}かした\R{形}{かたち}（\R{赤}{あか}）で、\R{頂点}{ちょう|てん}がいちばん\R{高}{たか}い\R{点}{てん}です。
\end{minipage}
\end{reidai}

\begingroup\tcbset{equal height group=bkthrjB}%
\noindent\begin{minipage}[t]{0.49\linewidth}
\begin{jibun}
\raggedright $y=3x^2$ について、\R{次}{つぎ}の\R{変化}{へん|か}の\R{割合}{わり|あい}を\R{求}{もと}めましょう。\par（1）\nobreak $x$ が2から4まで\par（2）\nobreak $x$ が2から $2+h$ まで
\bkjans{（1）\nobreak $\dfrac{48-12}{4-2}=\dfrac{36}{2}=18$（2）\nobreak $3(2+h)^2=12+12h+3h^2$。$\dfrac{12h+3h^2}{h}=12+3h$}
\end{jibun}
\end{minipage}\hfill
\begin{minipage}[t]{0.49\linewidth}
\begin{jibun}
\raggedright \R{次}{つぎ}の\R{放物線}{ほう|ぶつ|せん}の\R{頂点}{ちょう|てん}と\R{軸}{じく}を\R{答}{こた}え、\R{上}{うえ}に\R{開}{ひら}くか\R{下}{した}に\R{開}{ひら}くかを\R{答}{こた}えましょう。\par（1）\nobreak $y=(x-1)^2-5$\par（2）\nobreak $y=3(x+2)^2-1$
\bkjans{（1）\nobreak \R{頂点}{ちょう|てん} $(1,\ -5)$、\R{軸}{じく} $x=1$、\R{上}{うえ}に\R{開}{ひら}く（2）\nobreak \R{頂点}{ちょう|てん} $(-2,\ -1)$、\R{軸}{じく} $x=-2$、\R{上}{うえ}に\R{開}{ひら}く}
\end{jibun}
\end{minipage}\par
\endgroup

\begin{reidai}[\R{平方完成}{へい|ほう|かん|せい}して\R{頂点}{ちょう|てん}を\R{求}{もと}める]
\R{次}{つぎ}の\R{放物線}{ほう|ぶつ|せん}の\R{頂点}{ちょう|てん}を\R{求}{もと}めましょう。\par（1）\nobreak $y=x^2-4x+7$\quad（2）\nobreak $y=-5t^2+30t$\par
\kangaekata （1）$x$ の\R{係数}{けい|すう} $-4$ の\R{半分}{はん|ぶん}の $-2$ を\R{使}{つか}い、$(x-2)^2$ を\R{作}{つく}ります。（2）\R{先}{さき}に $t^2$ の\R{係数}{けい|すう} $-5$ でくくってから、かっこの\R{中}{なか}を\R{平方完成}{へい|ほう|かん|せい}します。\par
\shiki（1）\nobreak $(x-2)^2=x^2-4x+4$ なので、4をたして4をひく\par
\hspace*{4.3\zw}$x^2-4x+7=x^2-4x+4-4+7=(x-2)^2+3$\par
\hspace*{2.6\zw}（2）\nobreak $-5t^2+30t=-5(t^2-6t)$\quad $-5$ でくくる\par
\hspace*{4.3\zw}$t^2-6t=t^2-6t+9-9=(t-3)^2-9$\quad $-6$ の\R{半分}{はん|ぶん}は $-3$\par
\hspace*{4.3\zw}$y=-5\{(t-3)^2-9\}=-5(t-3)^2+45$\quad $-5\times(-9)=45$\par
\kotae（1）\nobreak \R{頂点}{ちょう|てん} $(2,\ 3)$\quad（2）\nobreak \R{頂点}{ちょう|てん} $(3,\ 45)$\par
\noindent\begin{minipage}[t]{0.48\linewidth}\centering
\begin{tikzpicture}
\begin{axis}[bkplot, width=52mm, height=44mm, xmin=-0.6, xmax=4.8, ymin=-0.8, ymax=8.4, xtick={1,2,3,4}, ytick={3},
  xlabel={$x$}, ylabel={$y$}]
  \addplot[bkcurve, domain=-0.3:4.3] {(x-2)^2+3};
  \addplot[only marks, mark=*, mark size=2pt, bkRed] coordinates {(2,3)};
  \addplot[only marks, mark=*, mark size=1.5pt, bkInk] coordinates {(0,7) (4,7)};
\end{axis}
\end{tikzpicture}
\bkfigcaption{（1）$y=(x-2)^2+3$。$x=0$ で7（もとの\R{式}{しき}の7と\R{合}{あ}う）。}
\end{minipage}\hfill
\begin{minipage}[t]{0.48\linewidth}\centering
\begin{tikzpicture}
\begin{axis}[bkplot, width=52mm, height=44mm, xmin=0, xmax=6.8, ymin=0, ymax=52, xtick={1,...,6}, ytick={25,45},
  xlabel={$t$}, ylabel={$y$}]
  \addplot[bkcurve, domain=0:6] {-5*(x-3)^2+45};
  \addplot[only marks, mark=*, mark size=2pt, bkRed] coordinates {(3,45)};
\end{axis}
\end{tikzpicture}
\bkfigcaption{（2）$y=-5(t-3)^2+45$。\qty{30}{m/s} で\R{真上}{ま|うえ}に\R{投}{な}げたボールは、3\R{秒後}{びょう|ご}に\R{最高点}{さい|こう|てん} \qty{45}{m} に\R{届}{とど}く。}
\end{minipage}
\end{reidai}

\begin{reidai}[\R{範囲}{はん|い}があるときの\R{最大}{さい|だい}と\R{最小}{さい|しょう}]
\R{次}{つぎ}の\R{関数}{かん|すう}の\R{最大値}{さい|だい|ち}と\R{最小値}{さい|しょう|ち}を\R{求}{もと}めましょう。\par
（1）\nobreak $y=x^2$（$1\leqq x\leqq3$）\par（2）\nobreak $y=x^2$（$-1\leqq x\leqq2$）\par（3）\nobreak $y=-5(t-2)^2+20$（$0\leqq t\leqq3$）\par
\kangaekata \R{頂点}{ちょう|てん}が\R{範囲}{はん|い}の\R{中}{なか}にあるかを\R{見}{み}てから、はしの\R{値}{あたい}と\R{頂点}{ちょう|てん}の\R{値}{あたい}をくらべます。\par
\shiki（1）\nobreak \R{頂点}{ちょう|てん} $x=0$ は\R{範囲}{はん|い}の\R{外}{そと}。$x=1$ で1、$x=3$ で9\par
\hspace*{2.6\zw}（2）\nobreak \R{頂点}{ちょう|てん} $x=0$ は\R{範囲}{はん|い}の\R{中}{なか}で $y=0$。$x=-1$ で1、$x=2$ で4\par
\hspace*{2.6\zw}（3）\nobreak \R{頂点}{ちょう|てん} $t=2$ は\R{範囲}{はん|い}の\R{中}{なか}で $y=20$。\par
\hspace*{4.3\zw}$t=0$ で $-5\times4+20=0$、$t=3$ で $-5\times1+20=15$\par
\kotae（1）\nobreak \R{最大値}{さい|だい|ち}9（$x=3$）、\R{最小値}{さい|しょう|ち}1（$x=1$）\par
\hspace*{2.6\zw}（2）\nobreak \R{最大値}{さい|だい|ち}4（$x=2$）、\R{最小値}{さい|しょう|ち}0（$x=0$）\par
\hspace*{2.6\zw}（3）\nobreak \R{最大値}{さい|だい|ち}20（$t=2$）、\R{最小値}{さい|しょう|ち}0（$t=0$）
\end{reidai}

\begingroup\tcbset{equal height group=bkthrjC}%
\noindent\begin{minipage}[t]{0.49\linewidth}
\begin{jibun}
\raggedright \R{次}{つぎ}の\R{放物線}{ほう|ぶつ|せん}の\R{頂点}{ちょう|てん}を\R{求}{もと}めましょう。\par（1）\nobreak $y=x^2+2x-3$\par（2）\nobreak $y=-5t^2+10t$
\bkjans{（1）\nobreak $(x+1)^2-1-3=(x+1)^2-4$ で\R{頂点}{ちょう|てん} $(-1,\ -4)$（2）\nobreak $-5(t^2-2t)=-5\{(t-1)^2-1\}=-5(t-1)^2+5$ で\R{頂点}{ちょう|てん} $(1,\ 5)$}
\end{jibun}
\end{minipage}\hfill
\begin{minipage}[t]{0.49\linewidth}
\begin{jibun}
\raggedright \R{次}{つぎ}の\R{関数}{かん|すう}の\R{最大値}{さい|だい|ち}と\R{最小値}{さい|しょう|ち}を\R{求}{もと}めましょう。\par
（1）\nobreak\mbox{$y=2x^2$（$1\leqq x\leqq2$）}\par（2）\nobreak\mbox{$y=2x^2$（$-3\leqq x\leqq1$）}\par（3）\nobreak $y=-5(t-1)^2+5$\par\hspace*{2\zw}（$0\leqq t\leqq1.5$）
\bkjans{（1）\nobreak \R{最大値}{さい|だい|ち}8（$x=2$）、\R{最小値}{さい|しょう|ち}2（$x=1$）（2）\nobreak \R{最大値}{さい|だい|ち}18（$x=-3$）、\R{最小値}{さい|しょう|ち}0（$x=0$）（3）\nobreak \R{最大値}{さい|だい|ち}5（$t=1$）、\R{最小値}{さい|しょう|ち}0（$t=0$）}
\end{jibun}
\end{minipage}\par
\endgroup

\begin{reidai}[\R{横}{よこ}に\R{投}{な}げたボールの\R{通}{とお}り\R{道}{みち}]
\R{高}{たか}さ \qty{1.25}{m} の\R{台}{だい}のはしから、ボールを \qty{5}{m/s} の\R{速}{はや}さで\R{真横}{ま|よこ}に\R{投}{な}げ\R{出}{だ}しました。\R{横}{よこ}に $x$〔m〕\R{進}{すす}んだときに、\R{台}{だい}の\R{高}{たか}さから $y$〔m〕\R{下}{さ}がっているとします。（1）\nobreak $y$ を $x$ の\R{式}{しき}で\R{表}{あらわ}しましょう。（2）\nobreak 0.5\R{秒後}{びょう|ご}に\R{床}{ゆか}に\R{着}{つ}きました。\R{台}{だい}のはしから\R{何}{なん}m\R{先}{さき}に\R{着}{つ}きましたか。\R{式}{しき}で\R{確}{たし}かめましょう。\par
\kangaekata \R{横}{よこ}は\R{一定}{いっ|てい}の\R{速}{はや}さなので $x=5t$、たては $y=5t^2$。$x=5t$ から $t$ を $x$ の\R{式}{しき}で\R{表}{あらわ}し、$y=5t^2$ にかっこごと\R{入}{い}れます。\par
\shiki（1）\nobreak $x=5t$ の\R{両辺}{りょう|へん}を5でわって、$t=\dfrac{x}{5}$\par
\hspace*{4.3\zw}$y=5\left(\dfrac{x}{5}\right)^2=5\times\dfrac{x^2}{25}=\dfrac{1}{5}x^2=0.2x^2$\par
\hspace*{2.6\zw}（2）\nobreak $x=5\times0.5=2.5$\quad \R{確}{たし}かめ：$0.2\times2.5^2=0.2\times6.25=1.25$\par
\kotae（1）\nobreak $y=0.2x^2$\quad（2）\nobreak \qty{2.5}{m} \R{先}{さき}（\R{下}{さ}がった\R{長}{なが}さが\R{台}{だい}の\R{高}{たか}さ \qty{1.25}{m} と\R{合}{あ}う）
\end{reidai}

\begin{jibun}
\R{高}{たか}い\R{台}{だい}のはしから、ボールを \qty{10}{m/s} の\R{速}{はや}さで\R{真横}{ま|よこ}に\R{投}{な}げ\R{出}{だ}しました。\R{横}{よこ}に $x$〔m〕\R{進}{すす}んだときに $y$〔m〕\R{下}{さ}がっているとして、（1）\nobreak $y$ を $x$ の\R{式}{しき}で\R{表}{あらわ}しましょう。（2）\nobreak 0.4\R{秒後}{びょう|ご}には、\R{横}{よこ}に\R{何}{なん}m\R{進}{すす}み、\R{何}{なん}m\R{下}{さ}がっていますか。（1）の\R{式}{しき}でも\R{確}{たし}かめましょう。
\bkjans{（1）\nobreak $y=5\left(\dfrac{x}{10}\right)^2=0.05x^2$（2）\nobreak \qty{4}{m} \R{進}{すす}み、\qty{0.8}{m} \R{下}{さ}がる。$0.05\times4^2=0.8$ で\R{合}{あ}う}
\end{jibun}

%% ---------------- まちがえやすいところ・物理工学 ----------------
\Needspace{14\baselineskip}
\bksection{tsumazuki}{まちがえやすいところをチェック}\label{bk:3-9:tsuma}

\begin{tsumazuki}{2\R{乗}{じょう}に\R{比例}{ひ|れい}でも「2\R{倍}{ばい}で2\R{倍}{ばい}」？}
\bkng{$y=5x^2$ で $x$ を2\R{倍}{ばい}にすると、$y$ も2\R{倍}{ばい}}
\bkok{$x$ を2\R{倍}{ばい}にすると、$y$ は $2\times2=4$ \R{倍}{ばい}}
$x$ が2\R{回}{かい}かけられているので、2\R{倍}{ばい}も2\R{回}{かい}かかります。$x=1$ で5、$x=2$ で20 です。
\end{tsumazuki}

\begin{tsumazuki}{マイナスの $x$ の2\R{乗}{じょう}}
\bkng{$y=x^2$ で $x=-3$ のとき、$y=-9$}
\bkok{$y=(-3)^2=(-3)\times(-3)=9$}
マイナス×マイナスはプラスです（2-2）。だから\R{放物線}{ほう|ぶつ|せん}は\R{左右対称}{さ|ゆう|たい|しょう}になります。
\end{tsumazuki}

\begin{tsumazuki}{\R{頂点}{ちょう|てん}の\R{符号}{ふ|ごう}を\R{逆}{ぎゃく}にする}
\bkng{$y=(x-2)^2+1$ の\R{頂点}{ちょう|てん}は $(-2,\ 1)$}
\bkok{$x-2=0$ になる $x=2$ が\R{頂点}{ちょう|てん}なので、$(2,\ 1)$}
「かっこの\R{中}{なか}が0になる $x$」が\R{頂点}{ちょう|てん}の $x$ \R{座標}{ざ|ひょう}です。$y=(x+2)^2+1$ なら $(-2,\ 1)$ です。
\end{tsumazuki}

\begin{tsumazuki}{\R{範囲}{はん|い}の\R{最小}{さい|しょう}を、はしの\R{値}{あたい}にしてしまう}
\bkng{$y=x^2$（$-1\leqq x\leqq2$）の\R{最小値}{さい|しょう|ち}は、はしの $x=-1$ のときの1}
\bkok{\R{頂点}{ちょう|てん} $x=0$ が\R{範囲}{はん|い}の\R{中}{なか}にあるので、\R{最小値}{さい|しょう|ち}は0}
\R{頂点}{ちょう|てん}が\R{範囲}{はん|い}の\R{中}{なか}にあるかを、\R{先}{さき}に\R{確}{たし}かめましょう。グラフをかくと、まちがえません。
\end{tsumazuki}

\begin{tsumazuki}{\R{変化}{へん|か}の\R{割合}{わり|あい}を $y\div x$ で\R{計算}{けい|さん}する}
$y=2x^2$ で、$x$ が1から3まで\R{変}{か}わるときの\R{変化}{へん|か}の\R{割合}{わり|あい}は？
\bkng{$x=3$ のとき $y=18$ なので、$18\div3=6$}
\bkok{$\dfrac{18-2}{3-1}=\dfrac{16}{2}=8$}
\R{変化}{へん|か}の\R{割合}{わり|あい}は「\R{変化量}{へん|か|りょう}」どうしのわり\R{算}{ざん}です（3-5）。$y$ も $x$ も、あと－まえをしてからわります。
\end{tsumazuki}

\begin{tsumazuki}{\R{平方完成}{へい|ほう|かん|せい}で、くくった\R{数}{かず}をかけ\R{忘}{わす}れる}
\bkng{$-5(t^2-4t)=-5\{(t-2)^2-4\}=-5(t-2)^2-4$}
\bkok{$-5\{(t-2)^2-4\}=-5(t-2)^2+20$}
$-5$ は、かっこの\R{中}{なか}の $-4$ にもかかります。$-5\times(-4)=20$ です。$t=0$ を\R{入}{い}れて、もとの\R{式}{しき}（0）と\R{合}{あ}うか\R{確}{たし}かめましょう。
\end{tsumazuki}

\begin{butsuri}
\noindent\begin{minipage}[c]{0.6\linewidth}
\par\addvspace{3pt}\textbf{パラボラアンテナ}\par\noindent \R{衛星}{えい|せい}\R{放送}{ほう|そう}のアンテナの\R{皿}{さら}は、\R{放物線}{ほう|ぶつ|せん}を\R{軸}{じく}のまわりに\R{回}{かい}した\R{形}{かたち}をしています。\R{約}{やく} 3\R{万}{まん}6000 km はなれた\R{人工衛星}{じん|こう|えい|せい}からの\R{電波}{でん|ぱ}は、ほとんど\R{平行}{へい|こう}にやってきます。\R{軸}{じく}に\R{平行}{へい|こう}に\R{入}{はい}ってきた\R{電波}{でん|ぱ}は、\R{皿}{さら}で\R{反射}{はん|しゃ}すると、すべて1つの\R{点}{てん}（\textbf{\R{焦点}{しょう|てん}}）に\R{集}{あつ}まります。そこに\R{受信}{じゅ|しん}する\R{部品}{ぶ|ひん}を\R{置}{お}くと、\R{弱}{よわ}い\R{電波}{でん|ぱ}を\R{強}{つよ}めて\R{受}{う}けとれるのです。パラボラは、\R{放物線}{ほう|ぶつ|せん}（parabola）のことです。
\end{minipage}\hfill
\begin{minipage}[c]{0.36\linewidth}\centering
\begin{tikzpicture}[x=0.85mm, y=0.85mm]
  % 放物線 x = y^2/(4f)、f=10。皿は y=-16〜16、焦点 (10,0)
  \draw[bkGray!90!black, line width=1.6pt, domain=-16:16, samples=40] plot ({\x*\x/40},\x);
  \foreach \yy in {-12,-6,6,12} {
    \pgfmathsetmacro\xx{\yy*\yy/40}
    \draw[bkBlue, line width=0.6pt] (38,\yy) -- (\xx+0.6,\yy);
    \draw[-{Stealth[length=2.5pt]}, bkBlue, line width=0.6pt] (\xx+0.6,\yy) -- ($(\xx+0.6,\yy)!0.92!(10,0)$);}
  \fill[bkRed] (10,0) circle (1.1);
  \node[font=\sffamily\scriptsize, text=bkRed!80!black, anchor=north] at (14,-2.5) {\R{焦点}{しょう|てん}};
  \node[font=\sffamily\scriptsize, text=bkBlue!80!black, anchor=south] at (30,13) {\R{電波}{でん|ぱ}};
\end{tikzpicture}
\end{minipage}

\vspace{2pt}
\par\addvspace{3pt}\textbf{\R{止}{と}まるまでの\R{距離}{きょ|り}は、\R{速}{はや}さの2\R{乗}{じょう}に\R{比例}{ひ|れい}}\par\noindent
\noindent\begin{minipage}[c]{0.6\linewidth}
ブレーキをかけると、\R{車}{くるま}は1\R{秒}{びょう}ごとに\R{同}{おな}じだけおそくなります（ぬれた\R{道路}{どう|ろ}で\R{約}{やく} \qty{5}{m/s} ずつ）。\qty{10}{m/s} なら2\R{秒}{びょう}、\qty{20}{m/s} なら4\R{秒}{びょう}で\R{止}{と}まります。\R{止}{と}まるまでの\R{距離}{きょ|り}は、$v$-$t$ グラフの\R{三角形}{さん|かく|けい}の\R{面積}{めん|せき}です（2-3）。この\R{三角形}{さん|かく|けい}は、\R{同}{おな}じ\R{形}{かたち}を2つ\R{合}{あ}わせると\R{長方形}{ちょう|ほう|けい}（\R{点線}{てん|せん}）になるので、\R{面積}{めん|せき}は\R{長方形}{ちょう|ほう|けい}（\R{底辺}{てい|へん}×\R{高}{たか}さ）の\R{半分}{はん|ぶん}です（4-5 でくわしく）。\R{速}{はや}さが2\R{倍}{ばい}なら、\R{底辺}{てい|へん}（\R{時間}{じ|かん}）も\R{高}{たか}さ（\R{速}{はや}さ）も2\R{倍}{ばい}で、\R{面積}{めん|せき}は4\R{倍}{ばい}です。\qty{10}{m/s}（\R{時速}{じ|そく} \qty{36}{km}）で \qty{10}{m} なら、\qty{20}{m/s}（\R{時速}{じ|そく} \qty{72}{km}）では \qty{40}{m} です。スピードの\R{出}{だ}しすぎが\R{危}{あぶ}ないのは、このためです。
\end{minipage}\hfill
\begin{minipage}[c]{0.36\linewidth}\centering
\begin{tikzpicture}
\begin{axis}[bkplot, width=38mm, height=38mm, xmin=0, xmax=4.9, ymin=0, ymax=24, xtick={2,4}, ytick={10,20},
  xlabel={$t$（s）}, ylabel={$v$（m/s）}]
  \addplot[bkOrange!70, line width=0.7pt, dash pattern=on 2pt off 1.5pt] coordinates {(0,20) (4,20) (4,0)};
  \addplot[fill=bkOrange!25, draw=bkOrange, line width=1pt] coordinates {(0,20) (4,0) (0,0)} -- cycle;
  \addplot[fill=bkBlue!25, draw=bkBlue, line width=1pt] coordinates {(0,10) (2,0) (0,0)} -- cycle;
  \node[font=\sffamily\scriptsize, text=bkBlue!80!black] at (axis cs:0.6,3) {10};
  \node[font=\sffamily\scriptsize, text=bkOrange!80!black] at (axis cs:1.9,7.5) {40};
\end{axis}
\end{tikzpicture}
\end{minipage}
\end{butsuri}

\begin{butsuri}
\par\addvspace{3pt}\textbf{\R{運動}{うん|どう}エネルギーも、\R{速}{はや}さの2\R{乗}{じょう}に\R{比例}{ひ|れい}}\par\noindent \R{動}{うご}いている\R{物}{もの}は、ほかの\R{物}{もの}をおしたり、こわしたりする\R{力}{ちから}をもっています。これを\textbf{\R{運動}{うん|どう}エネルギー}といい、\R{質量}{しつ|りょう} $m$〔kg〕、\R{速}{はや}さ $v$〔m/s〕の\R{物}{もの}では $\dfrac{1}{2}mv^2$〔J（ジュール）〕です。$\dfrac{1}{2}$ は、\R{上}{うえ}のブレーキの\R{三角形}{さん|かく|けい}の「÷2」と\R{同}{おな}じ\R{所}{ところ}から\R{出}{で}てきます（\R{物理}{ぶつ|り}シリーズでくわしく）。\R{速}{はや}さが2\R{倍}{ばい}なら、エネルギーは4\R{倍}{ばい}。だから、\R{車}{くるま}の\R{衝突}{しょう|とつ}は、\R{速}{はや}さが2\R{倍}{ばい}になると、こわれ\R{方}{かた}がずっと\R{大}{おお}きくなります。

\par\vspace{6pt}\noindent\begin{minipage}[c]{0.58\linewidth}
\textbf{\R{噴水}{ふん|すい}の\R{水}{みず}}\par\noindent \R{噴水}{ふん|すい}のノズルから、ななめ\R{上}{うえ}にふき\R{出}{だ}した\R{水}{みず}の1つぶ1つぶは、\R{横}{よこ}には\R{一定}{いっ|てい}の\R{速}{はや}さで、たてには\R{投}{な}げ\R{上}{あ}げと\R{同}{おな}じに\R{動}{うご}きます。だから、\R{水}{みず}の\R{線}{せん}は\R{放物線}{ほう|ぶつ|せん}をえがきます。ノズルの\R{向}{む}きを\R{変}{か}えると、\R{放物線}{ほう|ぶつ|せん}の\R{高}{たか}さや\R{遠}{とお}くまで\R{届}{とど}く\R{長}{なが}さが\R{変}{か}わります。\R{同}{おな}じ\R{速}{はや}さでふき\R{出}{だ}すとき、いちばん\R{遠}{とお}くまで\R{届}{とど}くのは、ななめ45°に\R{向}{む}けたときです（\R{空気}{くう|き}の\R{抵抗}{てい|こう}がないとき）。
\end{minipage}\hfill
\begin{minipage}[c]{0.38\linewidth}\centering
\begin{tikzpicture}[x=0.85mm, y=0.85mm]
  % 速さ 10 m/s、角度 30°・45°・60°。1 m を 4 mm でかいた（g=10）
  \fill[bkBlue!15] (-4,-3) rectangle (46,0);
  \draw[bkBlue!70!black, line width=0.6pt] (-4,0) -- (46,0);
  \foreach \ang/\col in {30/bkGreen!70!black, 45/bkBlue, 60/bkRed} {
    \pgfmathsetmacro\R{10*10*sin(2*\ang)/10}
    \draw[\col, line width=1.2pt, domain=0:\R, samples=40] plot ({4*\x},{4*(\x*tan(\ang)-10*\x*\x/(2*100*cos(\ang)*cos(\ang)))});}
  \fill[bkGray!80] (-1.5,0) rectangle (1.5,1.5);
\end{tikzpicture}
\bkfigcaption{\R{赤}{あか}60°、\R{青}{あお}45°、\R{緑}{みどり}30°。}
\end{minipage}
\end{butsuri}

\begin{column}{つり\R{橋}{はし}のケーブルは、\R{放物線}{ほう|ぶつ|せん}}
\noindent\begin{minipage}[c]{0.56\linewidth}
つり\R{橋}{はし}では、\R{太}{ふと}いケーブルから、たくさんのつり\R{下}{さ}げ\R{用}{よう}のロープで、\R{道路}{どう|ろ}（\R{橋}{はし}げた）をつり\R{下}{さ}げています。\R{道路}{どう|ろ}の\R{重}{おも}さが、\R{横}{よこ}の\R{長}{なが}さあたり\R{同}{おな}じだけかかると、ケーブルはちょうど\R{放物線}{ほう|ぶつ|せん}の\R{形}{かたち}になります。\R{設計}{せっ|けい}する\R{人}{ひと}は、この\R{放物線}{ほう|ぶつ|せん}の\R{式}{しき}を\R{使}{つか}って、ケーブルの\R{長}{なが}さや、かかる\R{力}{ちから}を\R{計算}{けい|さん}します。

ケーブルだけをぶら\R{下}{さ}げたときの\R{形}{かたち}（くさりの\R{形}{かたち}）は、\R{放物線}{ほう|ぶつ|せん}によく\R{似}{に}ていますが、\R{少}{すこ}しちがう\R{曲線}{きょく|せん}になります。これは\R{第}{だい}5\R{巻}{かん}の\R{先}{さき}で\R{学}{まな}ぶ\R{曲線}{きょく|せん}です。
\end{minipage}\hfill
\begin{minipage}[c]{0.4\linewidth}\centering
\begin{tikzpicture}[x=0.8mm, y=0.8mm]
  \fill[bkBlue!12] (-2,-6) rectangle (52,-2);
  \fill[bkGray!60] (2,-6) rectangle (5,22);
  \fill[bkGray!60] (45,-6) rectangle (48,22);
  \draw[bkInk, line width=1.1pt, domain=3.5:46.5, samples=40] plot (\x,{4+16*((\x-25)/21.5)^2});
  \foreach \xx in {7.5,11.5,...,43.5} {\draw[bkInk!60, line width=0.4pt] (\xx,{4+16*((\xx-25)/21.5)^2}) -- (\xx,1.2);}
  \fill[bkOrange!70] (0,0) rectangle (50,1.6);
\end{tikzpicture}
\end{minipage}
\end{column}

\begin{column}{スマホで、\R{放物線}{ほう|ぶつ|せん}をつかまえよう}
\noindent\begin{minipage}[c]{0.6\linewidth}
\R{家族}{か|ぞく}や\R{友達}{とも|だち}に、ボールを\R{山}{やま}なりに\R{投}{な}げてもらい、スマホで\R{横}{よこ}から\R{動画}{どう|が}をとります。\R{動画}{どう|が}をコマ\R{送}{おく}りにして、\R{同}{おな}じ\R{時間}{じ|かん}ごとのボールの\R{位置}{い|ち}を、\R{画面}{が|めん}の\R{上}{うえ}で\R{記録}{き|ろく}します（\R{画面}{が|めん}の\R{座標}{ざ|ひょう}は、$y$ が\R{下向}{した|む}きのことが\R{多}{おお}いので\R{注意}{ちゅう|い}。3-2）。

いちばん\R{高}{たか}い\R{所}{ところ}を\R{頂点}{ちょう|てん}にして、\R{横}{よこ} $x$ と、\R{頂点}{ちょう|てん}から\R{下}{さ}がった\R{長}{なが}さ $y$ を\R{表}{ひょう}にし、\R{表計算}{ひょう|けい|さん}で $y\div x^2$ を\R{計算}{けい|さん}すると、どの\R{点}{てん}でもだいたい\R{同}{おな}じ\R{数}{かず}になります。ボールの\R{通}{とお}り\R{道}{みち}が、$y=ax^2$ の\R{放物線}{ほう|ぶつ|せん}だとわかります。ずれた\R{点}{てん}に、どの\R{放物線}{ほう|ぶつ|せん}がいちばん\R{合}{あ}うかを\R{決}{き}める\R{方法}{ほう|ほう}は、3-16 で\R{学}{まな}びます。
\end{minipage}\hfill
\begin{minipage}[c]{0.36\linewidth}\centering
\begin{tikzpicture}[x=1mm, y=1mm]
  \filldraw[fill=bkGray!20, draw=bkInk!80, line width=0.8pt, rounded corners=2mm] (0,0) rectangle (44,28);
  \draw[bkBlue!60, line width=0.6pt, dash pattern=on 1.5pt off 1.2pt, domain=4:40, samples=40] plot (\x,{22-0.05*(\x-22)*(\x-22)});
  \foreach \xx in {6,10,...,38} {\bkthrjball{(\xx,{22-0.05*(\xx-22)*(\xx-22)})}{1.2}}
  \fill[bkInk!80] (44.8,11) rectangle (46,17);
\end{tikzpicture}
\end{minipage}
\end{column}

%% ---------------- 章末問題 ----------------
\bkendproblems
\bkproblemgroup{\R{基本}{き|ほん}}{2}

\bkproblem $y$ は $x$ の2\R{乗}{じょう}に\R{比例}{ひ|れい}し、$x=3$ のとき $y=27$ です。$y$ を $x$ の\R{式}{しき}で\R{表}{あらわ}し、$x=-2$ のときの $y$ を\R{求}{もと}めましょう。
\bkpans{$a=\dfrac{27}{9}=3$ で $y=3x^2$。$x=-2$ のとき $y=3\times4=12$}

\bkproblem \R{次}{つぎ}の\R{関数}{かん|すう}について、$x=-2$、$-1$、0、1、2 のときの $y$ を\R{求}{もと}めて\R{表}{ひょう}にし、グラフをノートにかきましょう。\par\hspace*{1\zw}（1）\nobreak $y=\dfrac{1}{4}x^2$\quad（2）\nobreak $y=-x^2$
\bkpans{（1）\nobreak 1、0.25、0、0.25、1。\R{上}{うえ}に\R{開}{ひら}いた\R{放物線}{ほう|ぶつ|せん}。（2）\nobreak $-4$、$-1$、0、$-1$、$-4$。\R{下}{した}に\R{開}{ひら}いた\R{放物線}{ほう|ぶつ|せん}。どちらも\R{頂点}{ちょう|てん}は\R{原点}{げん|てん}}

\bkproblem $y=4x^2$ について、\R{次}{つぎ}の\R{変化}{へん|か}の\R{割合}{わり|あい}を\R{求}{もと}めましょう。\par\hspace*{1\zw}（1）\nobreak $x$ が1から3まで\quad（2）\nobreak $x$ が1から $1+h$ まで
\bkpans{（1）\nobreak $\dfrac{36-4}{2}=16$（2）\nobreak $4(1+h)^2-4=8h+4h^2$。$8+4h$}

\bkproblem \R{次}{つぎ}の\R{放物線}{ほう|ぶつ|せん}の\R{頂点}{ちょう|てん}と\R{軸}{じく}を\R{答}{こた}え、\R{上}{うえ}に\R{開}{ひら}くか\R{下}{した}に\R{開}{ひら}くかを\R{答}{こた}えましょう。\par\hspace*{1\zw}（1）\nobreak $y=(x+4)^2-1$\quad（2）\nobreak $y=-3(x-2)^2+5$
\bkpans{（1）\nobreak \R{頂点}{ちょう|てん} $(-4,\ -1)$、\R{軸}{じく} $x=-4$、\R{上}{うえ}に\R{開}{ひら}く（2）\nobreak \R{頂点}{ちょう|てん} $(2,\ 5)$、\R{軸}{じく} $x=2$、\R{下}{した}に\R{開}{ひら}く}

\bkproblem \R{次}{つぎ}の\R{放物線}{ほう|ぶつ|せん}の\R{頂点}{ちょう|てん}を、\R{平方完成}{へい|ほう|かん|せい}して\R{求}{もと}めましょう。\par\hspace*{1\zw}（1）\nobreak $y=x^2-6x+5$\quad（2）\nobreak $y=-5t^2+40t$
\bkpans{（1）\nobreak $(x-3)^2-4$ で\R{頂点}{ちょう|てん} $(3,\ -4)$（2）\nobreak $-5(t-4)^2+80$ で\R{頂点}{ちょう|てん} $(4,\ 80)$}

\bkproblem $y=x^2$ の\R{最大値}{さい|だい|ち}と\R{最小値}{さい|しょう|ち}を\R{求}{もと}めましょう。\par\hspace*{1\zw}（1）\nobreak $2\leqq x\leqq4$\quad（2）\nobreak $-3\leqq x\leqq1$
\bkpans{（1）\nobreak \R{最大値}{さい|だい|ち}16（$x=4$）、\R{最小値}{さい|しょう|ち}4（$x=2$）（2）\nobreak \R{最大値}{さい|だい|ち}9（$x=-3$）、\R{最小値}{さい|しょう|ち}0（$x=0$）}

\bkproblemgroup{チャレンジ}{3}

\bkproblem \qty{30}{m/s} の\R{速}{はや}さで\R{真上}{ま|うえ}に\R{投}{な}げ\R{上}{あ}げたボールの、$t$ \R{秒後}{びょう|ご}の\R{高}{たか}さは $y=30t-5t^2$（m）です。（1）\nobreak \R{平方完成}{へい|ほう|かん|せい}して、ボールが\R{最}{もっと}も\R{高}{たか}くなる\R{時刻}{じ|こく}と\R{高}{たか}さを\R{求}{もと}めましょう。（2）\nobreak $1\leqq t\leqq4$ のあいだの、\R{高}{たか}さの\R{最大値}{さい|だい|ち}と\R{最小値}{さい|しょう|ち}を\R{求}{もと}めましょう。
\bkpans{（1）\nobreak $y=-5(t-3)^2+45$。3\R{秒後}{びょう|ご}に \qty{45}{m}（2）\nobreak \R{頂点}{ちょう|てん} $t=3$ は\R{範囲}{はん|い}の\R{中}{なか}で45。$t=1$ で25、$t=4$ で40。\R{最大値}{さい|だい|ち} \qty{45}{m}（$t=3$）、\R{最小値}{さい|しょう|ち} \qty{25}{m}（$t=1$）}

\bkproblem \R{高}{たか}さ \qty{0.8}{m} の\R{机}{つくえ}のはしから、ボールが \qty{1}{m/s} で\R{真横}{ま|よこ}に\R{飛}{と}び\R{出}{だ}しました。\R{横}{よこ}に $x$〔m〕\R{進}{すす}んだときに $y$〔m〕\R{下}{さ}がっているとして、（1）\nobreak $y$ を $x$ の\R{式}{しき}で\R{表}{あらわ}しましょう。（2）\nobreak 0.4\R{秒後}{びょう|ご}に\R{床}{ゆか}に\R{着}{つ}きました。\R{机}{つくえ}のはしから\R{何}{なん}m\R{先}{さき}ですか。（1）の\R{式}{しき}で、\R{下}{さ}がった\R{長}{なが}さが \qty{0.8}{m} になることを\R{確}{たし}かめましょう。
\bkpans{（1）\nobreak $t=x$ なので $y=5x^2$（2）\nobreak $x=1\times0.4=0.4$ で \qty{0.4}{m} \R{先}{さき}。$5\times0.4^2=5\times0.16=0.8$ で\R{合}{あ}う}

\bkproblem \R{落}{お}ちるボールの、\R{落}{お}ちはじめてから $x$ \R{秒間}{びょう|かん}に\R{落}{お}ちる\R{距離}{きょ|り}は $y=5x^2$（m）です。（1）\nobreak 2\R{秒}{びょう}から3\R{秒}{びょう}までの1\R{秒間}{びょう|かん}の\R{平均}{へい|きん}の\R{速}{はや}さ（\R{変化}{へん|か}の\R{割合}{わり|あい}）を\R{求}{もと}めましょう。（2）\nobreak 2\R{秒}{びょう}から $2+h$ \R{秒}{びょう}までの\R{変化}{へん|か}の\R{割合}{わり|あい}を\R{求}{もと}め、$h=0.1$ のときの\R{値}{あたい}を\R{計算}{けい|さん}しましょう。
\bkpans{（1）\nobreak $\dfrac{45-20}{3-2}=25$ で \qty{25}{m/s}（2）\nobreak $5(2+h)^2-20=20h+5h^2$。$20+5h$。$h=0.1$ で20.5}

\bkproblem ぬれた\R{道路}{どう|ろ}で、\R{車}{くるま}が\R{止}{と}まるまでにブレーキで\R{進}{すす}む\R{距離}{きょ|り}は、\R{速}{はや}さの2\R{乗}{じょう}に\R{比例}{ひ|れい}します。\R{速}{はや}さ \qty{10}{m/s} のとき \qty{10}{m} でした。（1）\nobreak \R{速}{はや}さを $v$〔m/s〕、\R{距離}{きょ|り}を $y$〔m〕として、$y$ を $v$ の\R{式}{しき}で\R{表}{あらわ}しましょう。（2）\nobreak \qty{30}{m/s}（\R{時速}{じ|そく} \qty{108}{km}）のときの\R{距離}{きょ|り}を\R{求}{もと}めましょう。\R{速}{はや}さが3\R{倍}{ばい}なので、\R{距離}{きょ|り}は\R{何倍}{なん|ばい}ですか。
\bkpans{（1）\nobreak $a=\dfrac{10}{10^2}=0.1$ で $y=0.1v^2$（2）\nobreak $0.1\times900=90$ で \qty{90}{m}。9\R{倍}{ばい}}

\begin{furikaeri}
\begin{checklist}
\item 2\R{乗}{じょう}に\R{比例}{ひ|れい}を\R{見}{み}つけ、$y=ax^2$ の\R{式}{しき}とグラフにできる。\bkfurref{bk:3-9:nijo}
\item $a$ とグラフの\R{形}{かたち}の\R{関係}{かん|けい}や、\R{変化}{へん|か}の\R{割合}{わり|あい}を\R{説明}{せつ|めい}できる。\bkfurref{bk:3-9:a}
\item \R{平方完成}{へい|ほう|かん|せい}で\R{頂点}{ちょう|てん}を、\R{範囲}{はん|い}から\R{最大}{さい|だい}・\R{最小}{さい|しょう}を\R{求}{もと}められる。\bkfurref{bk:3-9:heihou}
\end{checklist}
\end{furikaeri}

\bkfukidashi{delta}{\R{投}{な}げ\R{上}{あ}げたボールは2\R{秒後}{びょう|ご}に \qty{20}{m}。では、\qty{15}{m} の\R{高}{たか}さを\R{通}{とお}るのは\R{何秒後}{なん|びょう|ご}かな？　\R{次}{つぎ}の\R{章}{しょう}で、2\R{次方程式}{じ|ほう|てい|しき}を\R{解}{と}いて\R{求}{もと}めよう。}
